\chapter{Wyomee Works}

\begin{description}
\item\annotation{Director:} Wyomee Dank (Carlos Thompson)
\item\annotation{Producer:} Suno
\item\annotation{Lyrics:} Some songs and sections are composed by Carlos Thompson; others are AI-generated (using Suno, Claude, ChatGPT, Gemini, or Grok). See individual song notes for specific attribution.
\end{description}

\sepbar

\section{El corrido de Carlson}
\tag{ricochets}\tag{john-carlson}\tag{narcocorrido}

\subsection{Notes}
Based on characters from \emph{Ricochets} by Carlos Thompson.\\
Character John R. Carlson.\\
Lyrics writen by Gemini on a prompt by the author. Low interaction.

\subsection{Style}
corrido norteño / narcocorrido

\begin{lyrics}

\songpart{Verso 1}
Se las dio de muy astuto,\\
el gringo Carlson mentado,\\
pensó que firmaba un pacto\\
con el diablo de aliado,\\
pero jugó con fuego,\\
y el negocio le ha fallado.

\songpart{Verso 2}
Por un lío de faldas viejas,\\
al cartel quiso usar,\\
les vio la cara de tontos\\
y se puso a celebrar,\\
pero las deudas de sangre\\
siempre te van a cobrar.

\songpart{Coro}
¡Ay, gringo Carlson, qué caro pagaste!\\
La ley de la mafia no sabe perdonar.\\
No te cobraron con oro ni plata,\\
fue tu pequeña la que vio el final.\\
Hoy tu arrogancia se viste de luto,\\
por lo que nunca vas a recuperar.

\songpart{Verso 3}
Dejó la ardiente Miami,\\
donde la dicha perdió,\\
buscando esconder el alma\\
donde no alumbrara el sol.\\
Llegó a la fría Bogotá,\\
con el pecho hecho carbón.

\songpart{Verso 4}
Hoy se le ve en las cantinas,\\
ahogando el frío en alcohol,\\
buscando en cada botella\\
el peso de su error.\\
Pero ni todo el aguardiente\\
le quita ese sinsabor.

\annotation{Tránsito / Hablado}
\emph{(Con música suave y el acordeón llorando)}
> "Y es que en este negocio, gringo...
> la astucia se paga con sangre.
> Ahora gane lo que gane, ya lo perdió todo."

\songpart{Verso Final}
Ahí deambula el gringo Carlson,\\
fantasma entre la neblina,\\
Miami le dio el dinero,\\
Bogotá la cruz divina.

El que la hace la paga,\\
así es la ley de la esquina.

\end{lyrics}

\section{Retrogrades}
\tag*{2023}\tag{social-media-seed}\tag{country}

\subsection{Notes}
Lyrics written by a LLM, based on a tweet by Carlos Thompson.

\subsection{Style}
country, bass, banjo, fiddle, harmonica; female lead

\begin{lyrics}
Some look up to the midnight sky,\\
And watch the planets roll,\\
To track the backward path they tread,\\
Upon a cosmic scroll.

They see the shadow cross the moon,\\
A dance of light and space,\\
A rare and wondrous clockwork move,\\
Across the stellar face.

\songpart{chorus}
Some of us look at the sky and retrogrades make our minds.\\
For some those planets when retrogrades bring fateful winds

For them it is a clock of stone,\\
Of science and of light,\\
An mathematical display\\
That decorates the night.

But others look upon the stars,\\
And feel a sudden dread,\\
They read the retrogrades with fear,\\
And bow a worried head.

\songpart{chorus}
Some of us look at the sky and retrogrades make our minds.\\
For some those planets when retrogrades bring fateful winds

They think the shadow on the moon\\
Will twist the ties of fate,\\
And bring a season of regret,\\
Or open up a gate.

To them the sky is whispering\\
Of fortunes lost or won,\\
A warning written in the dark,\\
Before the rising sun.

\songpart{chorus}
Some of us look at the sky and retrogrades make our minds.\\
For some those planets when retrogrades bring fateful winds

\end{lyrics}

\section{Your Role}
\tag*{2023}\tag{original-idea}\tag{ai-generated}\tag{punk-rock}

\subsection{Notes}
Poem written by ChatGPT based on a text by Carlos Thompson, not originally intended as lyrics.

\subsection{Style}
punk rock

\begin{lyrics}
A comrade to share life's grand charade, -\\
of aging side by side, our bond displayed.

A partner for discourse on worlds unknown, -\\
fantasy's realm explored, minds overthrown.

A playful ally, flirty words we send, -\\
texts or chats, a playful blend.

An outlet found, desires to set free, -\\
no strings attached, a space just to be.

A cuddle buddy, warmth in nights so deep, -\\
movies' glow or slumber's sweep.

A collaborator, schemes we devise, -\\
to take the world, power will rise.

An explorer bold, our planet's space, -\\
remote places, the unknown embrace.

A sender sweet, kitties' cute delight, -\\
on screens they play, smiles ignite.

A guardian watchful, complacency's foe, -\\
keep me on track, where growth will flow.

Of these diverse roles, which to show, -\\
to fill with grace, where talents grow?

\end{lyrics}

\section{De igual pra igual}
\tag{thomas-entourage}\tag{camila-alexander}\tag{xumi}\tag{piseiro}

\subsection{Notes}
Based on the \emph{Thomas' Entourage} fictional universe by Carlos Thompson.\\
Character: Camila Eugenia Santos Alexander, Xumi.\\
Co-written between Carlos Thompson and ChatGPT. High interaction.

\subsection{Style}
piseiro, female voice

\begin{lyrics}

\songpart{Verso 1}
No bar depois do jogo\\
ela escutava mais do que falava,\\
sabia reconhecer no riso fácil\\
quem ficava e quem passava.

\songpart{Verso 2}
Aprendeu cedo que carinho\\
não precisa se prender,\\
que amizade é a parte funda\\
que faz o resto acontecer.

\songpart{Refrão}
Xumi nunca quis um dono,\\
nem viver por condição,\\
quem entrava na sua vida\\
entrava de igual pra igual.

\songpart{Verso 3}
Sport jogando madrugada,\\
Barcelona em luz neon,\\
ela dividia mesas, cidades,\\
voos entre Heathrow e Congonhas.

\songpart{Refrão}
Xumi nunca quis um dono,\\
nem viver por condição,\\
quem entrava na sua vida\\
entrava de igual pra igual.

\songpart{Ponte}
Primeiro vinha a confiança.\\
Depois, talvez, o calor.\\
Sem hierarquia escondida\\
vestida de amor.

\songpart{Verso 4}
No Eixample a noite desce\\
sobre copos e conversas,\\
e ela sabe que há afetos\\
que duram sem promessa.

\songpart{Refrão}
Xumi nunca quis um dono,\\
nem viver por condição,\\
quem entrava na sua vida\\
entrava de igual pra igual.

(drums solo)

\songpart{Refrão}
Xumi nunca quis um dono,\\
nem viver por condição,\\
quem entrava na sua vida\\
entrava de igual pra igual.

(Fade)

\end{lyrics}

\section{Open Hearts}
\tag{thomas-entourage}\tag{camila-alexander}\tag{xumi}\tag{punk}

\subsection{Notes}
Based on the \emph{Thomas' Entourage} fictional universe by Carlos Thompson.\\
Character: Camila Eugenia Santos Alexander, Xumi.\\
Co-written between Carlos Thompson and ChatGPT. High interaction.

\subsection{Style}
punk

\begin{lyrics}

\songpart{Verse 1}
She learned the maps by memory,\\
Heathrow lights and London rain,\\
São Luís beneath the summer heat,\\
Barcelona trains.

\songpart{Verse 2}
Some lovers spoke in future tense,\\
in houses, rings, and names,\\
But Xumi moved through crowded rooms\\
untouched by borrowed frames.

\songpart{Chorus}
She was never waiting\\
for another life to start,\\
Built her world from chosen cities,\\
football nights, and open hearts.

\songpart{Verse 3}
Sport Club playing past midnight,\\
the screen a distant glow,\\
A Barça crowd, a Birmingham score,\\
three homes she’d always know.

\songpart{Chorus}
She was never waiting\\
for another life to start,\\
Built her world from chosen cities,\\
football nights, and open hearts.

\songpart{Bridge}
A friend before a lover.\\
A door without a chain.\\
No promise spoken falsely\\
just to make someone remain.

\songpart{Verse 4}
Autumn settled on the Eixample,\\
warm light across the bars,\\
And Xumi walked through Barcelona\\
beneath electric stars.

\songpart{Chorus}
She was never waiting\\
for another life to start,\\
Built her world from chosen cities,\\
football nights, and open hearts.

(drums solo)

\songpart{Chorus}
She was never waiting\\
for another life to start,\\
Built her world from chosen cities,\\
football nights, and open hearts.

(Fade)

\end{lyrics}

\section{Food Truck Line}
\tag{encounters}\tag{ai-generated}\tag{original-idea}\tag{indie-pop}

\subsection{Notes}
Lyrics written by Suno, based on a summary by Gemini on fictional characters and settings by Carlos Thompson.

\subsection{Style}
Indie-pop ballad with a steady mid-tempo pulse and soft swung percussion; verse rides sparse guitar, brushed kick, and muted bass; pre-chorus opens with rising strings and handclap lift; chorus lands with warm harmony stacks and a simple chant hook. Lead vocal stays close-mic and intimate, with doubled choruses, brief ad-lib tails, vinyl crackle between phrases, and a final lift from wide pads and tom swells. Crisp, cinematic, and tenderly detailed., gentle, stylish, distinct, warm

\begin{lyrics}

\songpart{Verse 1}
Sun on the pavement,\\
paper bags in a row\\
You stand in your jacket,\\
like you already know\\
Black hair in a bun,\\
green coat over black\\
Crossbody at your hip,\\
no need to look back

\songpart{Verse 2}
I see the way you hold\\
that quiet in your face\\
Curvy like a curve\\
that never lost its grace\\
He’s two steps behind you,\\
gray at the side\\
White shirt, blue-green tie,\\
easy little smile

\songpart{Pre-Chorus}
No words between you\\
Still the day feels close\\
Same line, same hunger,\\
same joke nobody told

\songpart{Chorus}
Waiting in the line\\
Side by side\\
Same noon, same shine\\
Different lives\\
Waiting in the line\\
You and I
(Food truck sky)
(Food truck sky)

\songpart{Verse 3}
The truck door swings open,\\
steam cuts through the heat\\
Someone calls an order,\\
somebody misses seats\\
Your boots in the dust,\\
his shoes clean and neat\\
Still it feels familiar\\
when the moment meets

\songpart{Pre-Chorus}
No words between you\\
Still the day feels close\\
Same line, same hunger,\\
same joke nobody told

\songpart{Chorus}
Waiting in the line\\
Side by side\\
Same noon, same shine\\
Different lives\\
Waiting in the line\\
You and I
(Food truck sky)
(Food truck sky)

\songpart{Bridge}
Maybe that’s the secret\\
of a working day\\
Two strangers in the open\\
who don’t look away\\
A grin, a jacket,\\
a glance, then gone\\
But for one warm minute,\\
it all feels like home

\songpart{Chorus}
Waiting in the line\\
Side by side\\
Same noon, same shine\\
Different lives\\
Waiting in the line\\
You and I
(Food truck sky)
(Food truck sky)

\end{lyrics}

\section{Thursday Counter}
\tag{ai-generated}\tag{original-idea}\tag{indie-folk}

\subsection{Notes}
Lyrics written by Suno, based on a product by Gemini on fictional characters and settings by Carlos Thompson.

\subsection{Style}
Indie folk with brushed acoustic guitar, upright bass, and fingerpicked banjo; verse stays sparse and close-mic with room tone, pre-chorus opens into soft harmony stacks, chorus adds hushed gang vocals and a rising cello line. Gentle swung groove, pencil-tap percussion, reverse guitar swells into transitions, vinyl hiss and wooden creaks for ear candy. Warm, intimate, slightly dusty mix with a hand-made feel., folk, slow

\begin{lyrics}

\songpart{Verse 1}
This Thursday leans in close\\
Like it knows my name\\
I stand at the counter\\
Longer than I came\\
I am not thinking hard\\
Just hovering near\\
Which might be a kind of thought\\
I keep around here

\songpart{Pre-Chorus}
The glass, the spoon, the quiet\\
Hold still for a while\\
Outside, the day keeps moving\\
With its neutral smile\\
I run the day ahead\\
Over and over slow\\
Like a loose tooth tested\\
Just to see what knows

\songpart{Chorus}
Five conversations\\
Five small weather systems\\
Five conversations\\
I carry them with me\\
I know all of them\\
I don’t know them fully\\
That’s the shape of morning\\
That’s the shape of me

\songpart{Verse 2}
Monday left them folded\\
In a back room of my head\\
Each one opening a little\\
When I’m trying to get fed\\
One asks for patience\\
One asks for a sign\\
One only needs a yes\\
That I can’t give on time\\
And every face I’ve met\\
Has had that hidden door\\
Half-closed, half-lighted\\
Than I noticed before

\songpart{Pre-Chorus}
So I leave the plate\\
Leave the toast behind\\
Walk out with the day\\
Unfinished in my mind\\
The air takes me in\\
Like it has done all week\\
No welcome, no warning\\
Just the pavement under me

\songpart{Chorus}
Five conversations\\
Five small weather systems\\
Five conversations\\
I carry them with me\\
I know all of them\\
I don’t know them fully\\
That’s the shape of morning\\
That’s the shape of me

\songpart{Bridge}
Maybe that’s the mercy\\
Of a day like this\\
No grand confession\\
Just the weight of what you miss\\
We are mostly strangers\\
Even in the same room\\
Compost in the back of the mind\\
Blooming when it will bloom

\annotation{Final Chorus}\\
Five conversations\\
Five small weather systems\\
Five conversations\\
I carry them with me\\
I know all of them\\
I don’t know them fully\\
That’s the shape of morning\\
That’s the shape of me

\songpart{Outro}
The city takes me in\\
Without a word\\
I am one more person\\
On its quiet errand tour

\end{lyrics}

\section{Kimizu Mamba Vit}
\tag{ai-generated}\tag{original-idea}\tag{techno}

\subsection{Notes}
Lyrics written by Suno, based on a product by Gemini on fictional characters and settings by Carlos Thompson.

\subsection{Style}
Techno with new age shimmer at 132 BPM, four-on-the-floor pulse and crisp syncopated hi-hats; verse rides a tight bass grid with glassy plucks, pre-chorus opens into rising arps and filtered pads, chorus hits with a bright saw lead and chant doubles on the names. Bridge drops to hand percussion, breathy stacked harmonies, and a reversed swell. Clean, glossy, futuristic mix with deep sub focus., clean, classic, soulful, warm, complex, world, beautiful, techno, rhythmic, new age

\begin{lyrics}

\songpart{Verse 1}
Kimizu in the feed\\
one clean line\\
sharp hands, sharp mind\\
fix the maze in real time

She trims the noise\\
cuts the drag\\
one touch, one answer\\
that’s her flag

\songpart{Pre-Chorus}
Mischa Iova\\
same name, same face\\
no costume in the mirror\\
no borrowed grace

She says it plain\\
and it lands like stone\\
every room she walks in\\
she stands alone

\songpart{Chorus}
Kimizu, Kimizu\\
clean and fast\\
Kimizu, Kimizu\\
built to last\\
Mischa Iova\\
say it clear\\
Mischa Iova\\
nothing hidden here

\songpart{Verse 2}
Marianne smiles\\
tea cup slow\\
polite like winter\\
soft enough to show

But the floor shifts\\
when she leans in close\\
Mamba Vit\\
and the room all knows

She moves like a blade\\
under silk and light\\
one quiet blink\\
then she wins the fight

\songpart{Pre-Chorus}
Mamba Vit\\
sweet on the tongue\\
one cold snap\\
and the deal gets done

Velvet glove\\
iron spine\\
she plays it gentle\\
then draws the line

\songpart{Chorus}
Kimizu, Kimizu\\
clean and fast\\
Kimizu, Kimizu\\
built to last\\
Mischa Iova\\
say it clear\\
Mischa Iova\\
nothing hidden here

\songpart{Bridge}
Three women, three codes\\
one pulse, one city\\
one is the fix\\
one is the truth\\
one is the pretty and deadly

Kimizu cuts through\\
Mischa holds true\\
Marianne turns\\
and the whole room bends for you

\songpart{Chorus}
Kimizu, Kimizu\\
clean and fast\\
Kimizu, Kimizu\\
built to last\\
Mischa Iova\\
say it clear\\
Mischa Iova\\
nothing hidden here\\
Mamba Vit\\
soft then bite\\
Mamba Vit\\
owns the night

\end{lyrics}

\section{Alarm Cabinet}
\tag*{2024}\tag{social-media-seed}\tag{techno}

\subsection{Notes}
Lyrics written by Suno, based on thread in X (Twitter) by Carlos Thompson.


\subsection{Style}
Techno with low-pitch guttural male vocals, 132 BPM, heavy syncopated kick pulse, rattling hi-hats, thick sub-bass, and machine-room stabs; verse rides sparse drum grid and cold synth ticks, pre-chorus tightens with filtered noise and vocal distortions, chorus hits with doubled bass growls and gang shouts on the anchor line; add alarm bleeps, reversed metallic swells, and short delay throws on key words; mix is dark, glossy, and chest-rattling, low, techno

\begin{lyrics}

\songpart{Verse 1}
My phone screams red again\\
I swipe it blind\\
Three alarms, four alarms\\
Still I cut the line

You call to check my steps\\
I tell you I'm on track\\
I say "yeah, I'm getting there"\\
Then I never call you back

\songpart{Pre-Chorus}
I can look busy\\
I can look clean\\
Folders in a row\\
Like a perfect machine

But I fold in half\\
When nobody sees\\
Hide the delay\\
Wear the fake receipts

\songpart{Chorus}
I need guided hands\\
I need guided hands\\
I need guided hands\\
When I disappear (guided hands)

I need a room with rules\\
A face that stays\\
A place that knows my lies\\
And keeps me in the frame

\songpart{Verse 2}
I sit in bright white rooms\\
With a clipboard stare\\
I move the paper around\\
Like I meant to care

Tabs in a neat little line\\
Cursor parked and still\\
I fake the motion well\\
But I chase the will

I build little walls of proof\\
Just to feel okay\\
Then slip out through the crack\\
And waste another day

\songpart{Pre-Chorus}
I don't need a slogan\\
I don't need a speech\\
I need a held-closed door\\
And someone in reach

A group with names like mine\\
A guide who won't drift\\
Something steady enough\\
To catch my hands when I split

\songpart{Chorus}
I need guided hands\\
I need guided hands\\
I need guided hands\\
When I disappear (guided hands)

I need a room with rules\\
A face that stays\\
A place that knows my lies\\
And keeps me in the frame

\songpart{Bridge}
Maybe I'm not broken\\
Maybe I'm just lost\\
In a maze of yes and no\\
And a count of what it cost

So let me sit in truth\\
Let me say it plain\\
I keep building alibis\\
Then I wear their chains

\annotation{Final Chorus}\\
I need guided hands\\
I need guided hands\\
I need guided hands\\
When I disappear (guided hands)

I need a room with rules\\
A face that stays\\
A place that knows my lies\\
And keeps me in the frame

I need guided hands\\
I need guided hands\\
I need guided hands\\
And some way to stay (guided hands)

\end{lyrics}

\section{Callejón Azul}
\tag*{2026}\tag{instrumental}\tag{jazz-cumbia}

\subsection{Notes}
Prompt rewritten by Suno based on original input by Carlos Thompson.

\subsection{Prompt}
Jazz cumbia instrumental with bass and percussion locked to a cumbia pulse, but the harmonic language stays smoky and urban; keyboards and electric guitar carry a lyrical progression while a second guitar slips into free-flowing jazz improvisations. Verse-like sections keep the groove tight and restrained, pre-chorus lifts with brighter voicings and syncopated fills, then the main theme swells and briefly breaks into a more melancholy turnaround. Smooth instrumental swaps, conga pickups, brushed transitions, and a short guitar run before the final drop. Warm, spacious mix with a polished club-jazz sheen., jazz, tropical, melancholy, cumbia, uplifting, electric

\sepbar

\section{Peccato d'aprile}
\tag{thomas-entourage}\tag{sylvia-lai}\tag{italian-folk-pop}

\subsection{Notes}
Based on the \emph{Thomas' Entourage} fictional universe by Carlos Thompson.\\
Character: Sylvia Lai.\\
Co-written between Carlos Thompson and Claude. High interaction.

Suno rewrote the lyrics to strip words in Campidanese, English, or French.

\subsection{Style}
folk-pop italiano / cadenze sarde\\
tempo: moderato, ♩≈ 76\\
female voice\\
early 2020s

\begin{lyrics}

\songpart{Verse 1}
Carbonia al mattino odora di macchia e sale\\
mia nonna conta i giorni come grani di corallo\\
il mare non risponde, ma io conosco il suo linguaggio\\
una ragazza stretta dentro un’isola troppo small\\
Avevo sedici anni e un bagaglio già pesante\\
ma marzo mi chiamava con una voce differente

\songpart{Verse 2}
Un aereo, poi una città di pietra,\\
medievale e fredda sotto un cielo color cenere,\\
un appartamento stretto, una finestra sul silenzio,\\
e un uomo con gli occhi scuri che mi guardava crescere.\\
Due persone e un inverno che non voleva finire,\\
io imparavo il nome esatto di ogni cosa da sentire

\songpart{Pre-Chorus}
E non era paura, no —\\
era qualcosa che sapevo già,\\
scritto prima che nascessi\\
nella pietra di quest'isola.\\
Sa mama mia m'at nau: no.\\
Ma il sangue dice altro

\songpart{Chorus}
Peccato d'aprile, peccato mio —\\
ho scelto di crescere fuori dal recinto,\\
non per cadere ma per capire\\
dove finisce la bambina e dove inizio.\\
Peccato d'aprile, e lo rifaccio —\\
e quella notte ho bruciato il calendario\\
per diventare quello che già ero

\songpart{Verse 3}
Aprile diventava maggio, maggio diventava estate,\\
le strade vuote fuori, dentro tutto più vicino,\\
leggevo i suoi silenzi come mappe di un territorio\\
che nessuna scuola insegna, che nessun libro ha scritto.\\
La sera cuocevamo, litigavamo di niente,\\
i confini diventavano una domanda senza risposta

\songpart{Pre-Chorus}
E non era colpa, no —\\
era qualcosa che cercavo già,\\
non lui, ma quello spazio\\
dove smettevo di aspettare.\\
Sa mama mia m'at nau: no.\\
Ma il sangue dice altro

\songpart{Chorus}
Peccato d'aprile, peccato mio —\\
ho scelto di crescere fuori dal recinto,\\
non per cadere ma per capire\\
dove finisce la bambina e dove inizio.\\
Peccato d'aprile, e lo rifaccio —\\
e quella notte ho bruciato il calendario\\
per diventare quello che già ero

\songpart{Bridge}
Sa bidda faghet rumoure —\\
il paese fa rumore,\\
la nonna che non dorme,\\
le voci sopra il tavolo.\\
Deo so issa chi fui —\\
solo che adesso so\\
il peso esatto di un sì

\songpart{Verse 4}
La mattina dopo abbiamo parlato come adulti,\\
con quella chiarezza strana che viene dopo il buio,\\
lui disse quello che pensava, io dissi quello che volevo —\\
nessuno recitava, nessuno aveva paura.\\
Otto maggio, un nome nuovo — protégée, accordo, forma —\\
una ragazza che aveva scelto la sua storia

\songpart{Pre-Chorus}
E non era fine, no —\\
era qualcosa che iniziava già,\\
non lui, ma quello che\\
di me non sapevo ancora.\\
Sa mama mia m'at nau: no.\\
Io ho detto: lo so, mamma\\
(Lo so)

\songpart{Chorus}
Peccato d'aprile, peccato mio —\\
ho scelto di crescere fuori dal recinto,\\
non per cadere ma per capire\\
dove finisce la bambina e dove inizio.\\
Peccato d'aprile, e lo rifaccio —\\
e quella notte ho bruciato il calendario\\
per diventare quello che già ero

\songpart{Chorus}
Peccato d'aprile, peccato bello —\\
non chiedo perdono a nessun cielo,\\
ho preso il mondo con le mie mani\\
e l'ho piegato verso di me.\\
Peccato d'aprile — era il mio,\\
era aprile, ero io

\end{lyrics}

\section{More Honestly Than Hearts}
\tag{thomas-entourage}\tag{aurora-swansen}\tag{r-and-b}

\subsection{Notes}
Based on the \emph{Thomas' Entourage} fictional universe by Carlos Thompson.\\
Character: Aurora Swansen.\\
Co-written between Carlos Thompson and Claude. High interaction.

\subsection{Style}
R\&B fast-passed, piano, percussions, bass

\begin{lyrics}

\songpart{Verse 1}
\annotation{baritone vocals}\\
In Worcester winter kitchens cold,\\
she learned what vows conceal,\\
her father spoke of faithlessness\\
like something factual, real.

\songpart{Verse 2}
\annotation{contralto vocals}\\
At Clark she watched young lovers swear\\
forever through the smoke,\\
then leave before the seasons turned\\
with barely words to cloak.

\songpart{Pre-Chorus}
\annotation{contralto vocals}\\
So she stopped calling longing love,\\
stopped treating touch as fate,\\
found cleaner truths in borrowed rooms\\
than churches legislate.

\songpart{Chorus}
\annotation{baritone+contralto vocals}\\
Aurora says the body speaks\\
more honestly than hearts,\\
but every system built on rules\\
must tear that truth apart.

\songpart{Verse 3}
\annotation{baritone vocals}\\
Now Bogotá in mountain rain,\\
new language on her tongue,\\
an office tower dressed in glass,\\
a life half-buried, young.

\songpart{Pre-Chorus}
\annotation{baritone vocals}\\
He speaks to her in careful terms,\\
the things he cannot show,\\
and leaves his armor by the chair\\
before the sessions go.

\songpart{Chorus}
\annotation{baritone+contralto vocals}\\
Aurora says the body speaks\\
more honestly than hearts,\\
but every system built on rules\\
must tear that truth apart.

\songpart{Bridge}
\annotation{contralto vocals}\\
No candles.\\
No promises.\\
No “always” whispered soft.\\
Just measured breath, fluorescent light,\\
her shoes beside the loft.

And maybe that is colder love,\\
or maybe simply clear:\\
to hold somebody carefully\\
without demanding years.

\songpart{Verse 4}
\annotation{contralto vocals}\\
Yet sometimes after clients leave\\
and dawn invades the blinds,\\
the silence waits beside the bed\\
like something left behind.

\songpart{Pre-Chorus}
\annotation{contralto vocals}\\
She still believes what she became,\\
the framework holds its shape,\\
though certain nights the empty room\\
feels larger than escape.

\songpart{Chorus}
\annotation{baritone+contralto vocals}\\
Aurora says the body speaks\\
more honestly than hearts,\\
but every system built on rules\\
must tear that truth apart.

\songpart{Outro}
\annotation{contrato vocals}\\
No romance.\\
No repentance.\\
Just the city breathing blue.\\
A woman trained to heal the mind\\
through what the rules refuse.

\end{lyrics}

\section{Hidden Things' Keeper}
\tag{thomas-entourage}\tag{aurora-swansen}\tag{noir-jazz}

\subsection{Notes}
Based on the \emph{Thomas' Entourage} fictional universe by Carlos Thompson.\\
Character: Aurora Swansen.\\
Co-written between Carlos Thompson and Claude. High interaction.

\subsection{Style}
noir jazz, fast-passed, sax, piano, bass

\begin{lyrics}

\songpart{Verse 1}
\annotation{baritone vocals}\\
She learned to sit completely still\\
while strangers came apart,\\
to separate the spoken words\\
from what they meant at heart.

\songpart{Verse 2}
One man confessed a second life,\\
one swore his wife still knew,\\
a woman traced her wedding ring\\
while saying they were through.

\songpart{Pre-Chorus}
And every door that softly closed\\
left fragments in the air,\\
Aurora carried all of them\\
like perfume in her hair.

\songpart{Chorus}
Aurora keeps the hidden things\\
nobody else can hold,\\
the trembling truths behind clean smiles,\\
the quiet and untold.

\songpart{Verse 3}
Then came the man who spoke in plans,\\
in structures and design,\\
who chose his words like measured steps\\
across a narrow line.

\songpart{Pre-Chorus}
He never cried, he never begged,\\
just loosened piece by piece,\\
and left his silences with her\\
like artifacts for keeps.

\songpart{Chorus}
Aurora keeps the hidden things\\
nobody else can hold,\\
the trembling truths behind clean smiles,\\
the quiet and untold.

\songpart{Bridge}
No voices now.\\
No clients left.\\
The blinds admit the dawn.\\
The room still hums with absent lives\\
long after they are gone.

She wonders what remains of her\\
when no one needs confession,\\
when all the borrowed grief falls still\\
outside the bounds of session.

\songpart{Verse 4}
By morning Bogotá awakes,\\
the traffic starts below,\\
another guarded soul arrives\\
with something not to show.

\songpart{Pre-Chorus}
She straightens out the legal pad,\\
composure reassembled,\\
while somewhere underneath the calm\\
a smaller silence trembled.

\songpart{Chorus}
Aurora keeps the hidden things\\
nobody else can hold,\\
the trembling truths behind clean smiles,\\
the quiet and untold.

\songpart{Outro}
No witness\\
for the witness.\\
No archive for her own.\\
Just Aurora and a thousand lives\\
she carries home alone.

\end{lyrics}

\section{Official Plus One}
\tag{thomas-entourage}\tag{gabriele-wiater}\tag{soft-pop-ballad}

\subsection{Notes}
Based on the \emph{Thomas' Entourage} fictional universe by Carlos Thompson.\\
Character: Gabriele Wiater, Gabi.\\
Co-written between Carlos Thompson and ChatGPT. High interaction.

\subsection{Style}
Soft pop-ballad tempo with industrial techno instruments; synthesizers, brass percussions, bass-heavy, quieter and more intimate when vocals. female voice

\begin{lyrics}

\songpart{Verse 1}
New York beneath the second wave, the avenues half-lit,\\
Gabi watched the airport codes like promises unfit.

\songpart{Verse 2}
The elevator softly chimed, the doorframe filled with light,\\
and something teal was next to him she hadn’t planned that night.

\songpart{Pre-Chorus}
So she reached first for practicals, for boundaries and tone,\\
the kind of calm that lets you hide how quickly you were thrown.

\songpart{Chorus}
No titles. No promises.\\
Still, somehow, she stayed near.\\
Like being chosen carefully\\
without the usual words to hear.

\songpart{Verse 3}
Brooklyn wrapped in radiator heat, shared blankets, coffee steam,\\
a girl with teal-dyed hair who asked too much with quiet green.

\songpart{Pre-Chorus}
And somewhere in those winter nights the seams gave way under heat,\\
through casual hands and honest talk and truths that nearly slipped.

\songpart{Chorus}
No titles. No guarantees.\\
No neat romantic frame.\\
Just Thomas reaching out to her\\
without pretending they were named.

\songpart{Bridge}
“Girlfriend” sounded wrong somehow.\\
“Wife” was even worse.\\
So they built softer language there\\
between affection, joke, and curse.\\
Official plus one.\\
She laughed too hard.\\
Which probably meant yes.

\songpart{Verse 4}
The runway lights at Teterboro cut gold across the rain,\\
and Gabi climbed the aircraft stairs without asking what they became.

\songpart{Pre-Chorus}
The world outside was shutting down, but somewhere in between\\
she found a shape of closeness still that lived outside routine.

\songpart{Chorus}
No titles. No promises.\\
Still, somehow, it was real.\\
A place beside somebody else\\
without surrendering the wheel.

\songpart{Outro}
Second wave. Cabin lights.\\
The engines humming low.\\
And Gabi letting herself stay\\
where she once planned to go.

\end{lyrics}

\section{The Motion and the Spark}
\tag{thomas-entourage}\tag{sylvia-lai}\tag{dream-pop}

\subsection{Notes}
Based on the \emph{Thomas' Entourage} fictional universe by Carlos Thompson.\\
Character: Sylvia Lai.\\
Co-written between Carlos Thompson and ChatGPT. High interaction.

\subsection{Style}
dream-pop

\begin{lyrics}

\songpart{Verse 1}
She came to the Big Apple once because Thomas was coming, one suitcase and a borrowed coat—\\
and why not follow something?

\songpart{Verse 2}
The second wave left empty streets, the cafés closing soon,\\
a gatekeeper with purple eyes,\\
and one more lover withdrawn too.

\songpart{Pre-Chorus}
But somewhere in those muted nights past curfew, smoke, and blue,\\
the girl beside the apartment door became somebody new.

\songpart{Chorus}
Sylvia still a girl then, still learning her own shape,\\
through half-lit rooms and winter talks and mornings waking late.\\
Not rescued. Not completed.\\
Not somebody passing through,\\
just slowly learning how to stay inside a life not fully new.

\songpart{Verse 3}
One night across the kitchen table, through coffee, wine, and steam,\\
they spoke about the kinds of love that blur with what they mean.

\songpart{Verse 4}
And maybe she loved Thomas less as flesh and bone and man,\\
than as the proof a different world might somehow still exist for her.

\songpart{Pre-Chorus}
And Gabi’s careful purple eyes would soften when she spoke,\\
with just enough exposed beneath the irony and smoke.

\songpart{Chorus}
Sylvia still a girl then, still learning how to feel,\\
through conversations after dark that somehow sounded real.\\
Not waiting to be chosen, not orbiting on cue,\\
just finding parts of who she was reflected back in someone new.

\songpart{Bridge}
And then she understood at last: Thomas had never been a destination or a prize to somehow finally win.\\
He was the road that brought her there, the motion and the spark, but she herself had crossed the bridge that winter through the dark.

\songpart{Verse 5}
New York landmarks through second wave: the skyline, bridge, and train,\\
the grocery on Atlantic Ave, the silence after rain.

\songpart{Pre-Chorus}
On December twenty-third she packed without regret,\\
and left behind the girl who came not knowing herself yet.

\songpart{Chorus}
Sylvia not just a girl now, not drifting place to place,\\
but someone who could meet the world without erasing space.\\
Not perfect. Not finished.\\
Still uncertain what was true,\\
but carrying home a quieter kind of personhood she grew into.

\songpart{Outro}
A year of learning. A month of learning.\\
The city washed in blue.\\
And Sylvia smiling in the glass like somebody she knew.

\end{lyrics}

\section{Casa non è}
\tag{thomas-entourage}\tag{sylvia-lai}\tag{italian-pop}

\subsection{Notes}
Based on the \emph{Thomas' Entourage} fictional universe by Carlos Thompson.\\
Character: Sylvia Lai.\\
Co-written between Carlos Thompson and Claude. High interaction.

\subsection{Style}
pop italiano\\
tempo: mid-tempo, ♩≈ 90\\
female voice\\
early 2020s

\begin{lyrics}

\songpart{Verse 1}
A Carbonia si contano i giorni dal mare,\\
la nonna sa il nome di ogni vento che arriva,\\
mia madre mi aspetta ogni sera alla porta,\\
mio padre non chiede ma tutto controlla.\\
Appartenevo a quel posto nel modo in cui appartieni a una cosa che ami e che non hai scelto.

\songpart{Verse 2}
Poi c'è stata una città che non esiste sulle mappe —\\
medievale e fredda, chiusa dentro se stessa,\\
un appartamento stretto, una finestra sul silenzio,\\
due persone e un inverno che non finiva mai.\\
Appartenevo a quel tempo nel modo in cui appartieni a una cosa che è reale ma non può restare.

\songpart{Pre-Chorus}
E poi l'Atlantico sotto, gli aeroporti piccoli,\\
una voce inglese precisa che leggevo senza capire,\\
la prima notte da sola in una stanza tutta mia —\\
non sapevo ancora come si chiama questa cosa.

\songpart{Chorus}
Casa non è il posto dove sei nata,\\
casa non è neanche dove resti,\\
casa è forse quella soglia che attraversi\\
quando smetti di chiedere dove appartieni.\\
Sono di passaggio in ogni luogo che abito,\\
e sto imparando che forse va bene così.

\songpart{Verse 3}
La famiglia lo accetta senza fare domande,\\
Francine chiede della Sardegna con quella grazia\\
che entra nei posti senza sembrare che entri —\\
diplomazia vera, quella che non si insegna.\\
Lucía mi abbraccia all'arrivo e alla partenza,\\
e non so cosa sa, e non chiedo niente.

\songpart{Pre-Chorus}
Ho detto: questo appartamento è mio.\\
L'ho detto piano, poi l'ho detto ad alta voce.\\
Ho deciso il colore prima di trovare dove farlo —\\
quando decido una cosa è già decisa.

\songpart{Chorus}
Casa non è il posto dove sei nata,\\
casa non è neanche dove resti,\\
casa è forse quella soglia che attraversi\\
quando smetti di chiedere dove appartieni.\\
Sono di passaggio in ogni luogo che abito,\\
e sto imparando che forse va bene così.

\songpart{Bridge}
C'è qualcosa in quest'uomo che riconosco\\
e non voglio guardare troppo da vicino —\\
una certezza che sembra di conoscere già,\\
che arriva prima che finisca la frase.\\
Il mare non c'è.\\
La città è troppo alta.\\
Ma resto perché sto imparando qualcosa di me.

\songpart{Verse 4}
A ottobre hanno firmato — una pratica chiusa,\\
lui era a suo agio come chi ha aspettato a lungo.\\
A dicembre ho restituito quello che mi aveva dato\\
quella notte di aprile con la torta e il vino —\\
la voglia non si duplica, ma ci siamo avvicinati,\\
e questo è il suo linguaggio e adesso è anche il mio.

\songpart{Pre-Chorus}
New York è già nel piano, dicembre è già vicino,\\
il piano giusto per adesso è quello che ho —\\
non cerco un posto dove restare per sempre,\\
cerco la prossima soglia da attraversare.

\songpart{Chorus}
Casa non è il posto dove sei nata,\\
casa non è neanche dove resti,\\
casa è forse quella soglia che attraversi\\
quando smetti di chiedere dove appartieni.\\
Sono di passaggio in ogni luogo che abito,\\
e sto imparando che forse va bene così.

\songpart{Chorus}
Casa non è — l'ho smesso di cercarla,\\
casa non è — forse non esiste,\\
casa è questa voce che si muove\\
da un confine all'altro senza chiedere permesso.\\
Sono di passaggio e il passaggio è mio,\\
e va bene così — va bene così.

\end{lyrics}

\section{Qualcosa di vero}
\tag{thomas-entourage}\tag{sylvia-lai}\tag{italian-pop}

\subsection{Notes}
Based on the \emph{Thomas' Entourage} fictional universe by Carlos Thompson.\\
Character: Sylvia Lai.\\
Co-written between Carlos Thompson and Claude. High interaction.

\subsection{Style}
dream-pop italiano\\
tempo: ♩≈ 80\\
female vocals

\begin{lyrics}

\songpart{VERSE 1}
\annotation{female vocals}\\
Teterboro piccola, la città non scatta ancora,\\
Upper East Side grigia come Bogotá ma meno calda,\\
Thomas costruisce un muro che non era lì prima,\\
freddo, schivo, distante senza spiegazione.\\
Gabi mi guarda dall'altra parte della stanza\\
come se stesse calcolando qualcosa che già sa.

\songpart{VERSE 2}
Ho passato più giorni con lei che con lui —\\
non era quello che mi aspettavo da New York.\\
Non fa domande, Gabi — fa osservazioni,\\
e le osservazioni fanno le domande al posto suo.\\
In una frase mi ha detto quello che pensavo su Claire —\\
l'ho archiviato, e sono andata avanti.

\songpart{PRE-CHORUS}
Due parole in due lingue,\\
un confine su una cartina —\\
mi fa arrabbiare un poco,\\
mi fa tenerezza, mi fa ridere.

\songpart{CHORUS}
Thomas era la persona o la sua idea?\\
Imparavo ancora la differenza.\\
Lui si era fatto da parte per una riga su una mappa,\\
e qualcosa di più caldo riempiva la stanza.\\
Non abbandonata, non persa —\\
solo dentro qualcosa che non sapevo ancora chiamare.

\songpart{VERSE 3}
Il suo studio a Queens — un disordine organizzato,\\
tele ovunque, luce buona, lei che si muoveva\\
intorno a me come cercasse qualcosa\\
che non sapeva ancora cosa fosse.\\
Non sessuale, non giudicante — solo attenta.\\
È un'esperienza essere guardata in quel modo.

\songpart{PRE-CHORUS}
Ha detto una cosa che non voleva dire,\\
le ho asciugato una lacrima con la mano.\\
Ho capito che potevo — e ho scelto di no.\\
Non tutto quello che si può fare si deve fare.

\songpart{CHORUS}
Thomas era la persona o la sua idea?\\
Imparavo ancora la differenza.\\
Lui si era fatto da parte per una riga su una mappa,\\
e qualcosa di più caldo riempiva la stanza.\\
Non aspettando, non in orbita —\\
solo trovando qualcosa di vero in una città di vetro.

\songpart{BRIDGE}
Quello che è successo in mezzo non si riassume —\\
sedici anni a marzo, qualcos'altro adesso.\\
Thomas era la strada, non la destinazione,\\
il movimento, la scintilla, il primo varco.\\
Ma Gabi mi ha guardata senza costruire una storia —\\
solo quello che ero, lì, in quella luce.

\songpart{VERSE 4}
Il ventidue dicembre faccio i bagagli senza rimpianti,\\
Newark–Roma–Cagliari, Greca e Alessandro all'aeroporto con la faccia.\\
Thomas è tornato caldo — non nel modo di prima,\\
nel modo giusto, quello tranquillo che basta.\\
Gabi mi ha fatto un regalo che mi ha fatto ridere —\\
Gabi è la mia persona preferita di New York.

\songpart{PRE-CHORUS}
Sono partita che avevo sedici anni,\\
torno con i capelli teal e qualcosa in più —\\
non so ancora dargli un nome,\\
ma lo riconosco quando ci sono dentro.

\songpart{CHORUS}
Thomas era la persona o la sua idea?\\
Adesso so la differenza.\\
Lui si era fatto da parte per una riga su una mappa,\\
e qualcosa di più vero riempiva la stanza.\\
Non persa, non salvata —\\
solo qualcuno che si riconosce nello specchio.

\songpart{CHORUS \annotation{finale}}
Qualcosa di vero in una città di vetro,\\
qualcosa di mio in un anno che non si riassume.\\
La strada era lui — la persona sono io,\\
e Gabi mi ha insegnato come si guarda davvero.\\
Non finita, non arrivata —\\
ma qualcuno che adesso sa dove guardare.

\end{lyrics}

\section{Sa fiamma}
\tag{thomas-entourage}\tag{sylvia-lai}\tag{italian-folk-pop}\tag{campidanese}

\subsection{Notes}
Based on the \emph{Thomas' Entourage} fictional universe by Carlos Thompson.\\
Character: Sylvia Lai.\\
Co-written between Carlos Thompson and Claude. High interaction.

\subsection{Style}
Italian pop, Sardinian folk influences, slow dream-pop. A steel-string acoustic guitar plays a steady eighth-note strumming pattern with a bright, percussive attack. A clean electric bass guitar follows the root notes of the chord progression with a warm, rounded tone. A female vocalist sings in a conversational, mezzo-soprano with a reflective, warm, bittersweet, resolved production style. The tempo is 72 BPM in 4/4 time, in the key of G major. The arrangement is sparse, focusing on the interplay between the rhythmic acoustic guitar and the vocal melody.

\begin{lyrics}

\songpart{Intro}
\annotation{bright acoustic guitar strumming, warm electric bass}

\songpart{Verse 1}
\annotation{female vocals}\\
Greca e Alessandro ad Elmas con la faccia —\\
nove mesi trattenuti in pubblico.\\
Le domande erano tutte sui capelli.
\annotation{Spoken word}\\
Giusto. I capelli sono quello che riescono a vedere.

\songpart{Verse 2}
Roberto sottovoce, una domanda sola —\\
gliene ho dette quattro in campidanese.\\
Tre ore di silenzio, poi è venuto a scusarsi.
\annotation{Spoken word}\\
Alcune cose funzionano meglio in campidanese.

\songpart{Pre-Chorus — Tallinn}
\annotation{Whisper}\\
Sono partita a marzo per due settimane —\\
il mondo ha chiuso i cancelli mentre ero dentro.\\
Ho scelto il peccato con gli occhi aperti —\\
e adesso lo porto senza chiedere perdono.

\songpart{Chorus}
\annotation{Joyful}\\
Sono tornata dove tutto è cominciato,\\
la villa, il fuoco, le voci che conosco.\\
Non sono piú quella che è partita a marzo —\\
ma il Truncu brucia uguale, e deve essere così.
\annotation{Spoken word}\\
Sa fiamma est sa fiamma (est sa fiamma)\\
cambia chi la guarda, non lei. (non lei)

\songpart{Verse 3}
Tutti i cugini, i mariti, i bambini che corrono,\\
la nonna che conta le teste intorno al fuoco.\\
La fiamma prende, il legno scoppietta, siamo tutti lì.
\annotation{Spoken word}\\
È ogni anno uguale e ogni anno necessario.

\songpart{Pre-Chorus — Bogotá}
\annotation{Whisper}\\
C'era un appartamento che ho chiamato mio\\
prima ancora di sapere se restavo.\\
Ho deciso il colore senza sapere dove farlo —\\
quando decido una cosa è già decisa.

\songpart{Chorus}
\annotation{Joyful}\\
Sono tornata dove tutto è cominciato,\\
la villa, il fuoco, le voci che conosco.\\
Non sono piú quella che è partita a marzo —\\
ma il Truncu brucia uguale, e deve essere così.
\annotation{Spoken word}\\
Sa fiamma est sa fiamma (est sa fiamma)\\
cambia chi la guarda, non lei. (non lei)

\songpart{Bridge}
\annotation{Vulnerable}\\
Quello che è successo in mezzo non si riassume:\\
una città di vetro, uno studio a Queens,\\
qualcuno che mi ha guardata senza costruire una storia.\\
Thomas era la strada — io ero la destinazione.\\
Avevo sedici anni, qualcos'altro adesso —
\annotation{Spoken word}\\
non so ancora il nome, ma lo riconosco.

\songpart{Verse 4}
Le foto di Gabi aperte sullo schermo —\\
un occhio quasi dorato, l'altro sotto i capelli teal.\\
Si contano le lentiggini, i puntini sul naso.
\annotation{Spoken word}\\
A volte li odio. In quella foto sono esattamente dove devono essere.

\songpart{Pre-Chorus — New York}
\annotation{Whisper}\\
Non quella che aspettava qualcuno che arrivasse,\\
non quella che chiedeva dove apparteneva.\\
Questa ragazza nello schermo la conosco —\\
potrei innamorarmi di me stessa.

\songpart{Chorus}
\annotation{Joyful}\\
Sono tornata dove tutto è cominciato,\\
la villa, il fuoco, le voci che conosco.\\
Non sono piú quella che è partita a marzo —\\
ma il Truncu brucia uguale, e deve essere così.
\annotation{Spoken word}\\
Sa fiamma est sa fiamma (est sa fiamma)\\
cambia chi la guarda, non lei. (non lei)

\annotation{Final Chorus — Campidanese, sparse arrangement, voice forward}\\
Seu torraa in dói chi naschìu,\\
sa villa, su fogu, is boghes chi connòsciu.\\
No seu prus issa chi partìu in màrtu —\\
ma su Truncu brusat sèmpri, e depit èssiri accussì.
\annotation{Spoken word}\\
Sa fiamma est sa fiamma —\\
mùdat chini la tzerriat, no issa.

\songpart{Outro — instrumental fade, acoustic guitar, single piano line}

\end{lyrics}

\section{Vuelta a Bal Harbour}
\tag{ricochets}\tag{fabian-muller}\tag{sandra-roa}\tag{r-and-b}

\subsection{Notes}
Writen by Suno based on a draft scene for \emph{Ricochets} by Carlos Thompson.

\subsection{Style}
Fast-paced R\&B with rap bridges, syncopated mid-tempo bounce and tight drum programming; verse rides sleek bass and clipped keys, pre-chorus opens with filtered harmonies and snapped ad-libs, chorus lands on stacked vocal hooks and a short chantable refrain. Rap bridges turn conversational and tense. Close-mic lead vocal, doubled hook lines, airy backing harmonies, quick risers, vinyl crackle hits, glossy and punchy mix., r\&b, clean, rap, light

\begin{lyrics}

\songpart{Verse 1}
Sandra vuelve de la playa,\\
arena en los zapatos.\\
Bal Harbour le quedó bien,\\
pero el mar ya quedó atrás.

Fabián ya estaba adentro,\\
cara serena, reloj calmado.\\
"¿Sabías de Carolina?"\\
Y él soltó: "Sí, Matteo."

"Italosuizo. Buena gente."\\
"¿Y me lo contó a ti?"\\
"No realmente", dijo él,\\
"Christian fue quien me lo presentó."

\songpart{Pre-Chorus}
"¿Christian quién?"\\
"Otro de Aachen."\\
"¿Sigue brava?"\\
"¿Conmigo?"\\
Sí, y Sandra aprieta la boca.

"Nicolás ya regresó?"\\
"No está aquí."\\
"Está en Bal Harbour Shops,"\\
dijo él, casi riéndose.

\songpart{Chorus}
No es justo, no es justo\\
pero aquí seguimos

No es justo, no es justo\\
y tú me sostienes

Bal Harbour, Bal Harbour\\
todo se ve tan limpio\\
Bal Harbour, Bal Harbour\\
y por dentro, qué frío

\songpart{Verse 2}
El Blackberry de Fabián vibra,\\
Perez sube por el pasillo.\\
Sandra corre por su cuarto,\\
cambia el paso, cambia el brillo.

Sale en business casual,\\
muy Miami, muy Cartagena.\\
Ella no lo piensa mucho,\\
solo entra y no se frena.

Khamid se va al instante,\\
queda el aire bien medido.\\
Y Fabián la presenta:\\
"Armando Perez, venido."

\songpart{Pre-Chorus}
"Agente especial", dice él,\\
"from the DEA."\\
"Call me Armando", y la mesa\\
se queda sin respiración.

"Estoy on hold por una operación,\\
si me llaman, tengo que salir."\\
Sandra mira a Fabián,\\
y él le aprieta la mano allí.

\songpart{Chorus}
No es justo, no es justo\\
pero aquí seguimos

No es justo, no es justo\\
y tú me sostienes

Bal Harbour, Bal Harbour\\
todo se ve tan limpio\\
Bal Harbour, Bal Harbour\\
y por dentro, qué frío

\songpart{Bridge}
"Let me be clear from the start,"\\
dice Pérez, sin parpadear,\\
"no somos el IRS,\\
pero sí te vamos a mirar."

"Clean" sus cuentas en Estados,\\
en Bahamas, y en Colombia,\\
pero el resto de la historia\\
huele a red, huele a sombra.

"Inter-Flota no va sola,"\\
dice, "todo está mezclado."\\
Sandra alza la barbilla:\\
"Eso suena a un monstruo inventado."

\annotation{Rap Bridge}\\
No es ficción, es la cuenta completa,\\
es sistema, no pieza suelta.\\
Te costó sangre levantarla,\\
y hoy te la quieren cerrar toda.

Metástasis en la estructura,\\
cáncer fino en la factura.\\
"No es justo", dice Sandra al filo,\\
y Fabián responde: "Yo sé, mi vida."

\songpart{Chorus}
No es justo, no es justo\\
pero aquí seguimos

No es justo, no es justo\\
y tú me sostienes

Bal Harbour, Bal Harbour\\
todo se ve tan limpio\\
Bal Harbour, Bal Harbour\\
y por dentro, qué frío

\songpart{Outro}
"Not free and clear,"\\
dice él, y cae la noche.\\
Sandra mira a Fabián:\\
"Esto todavía no es un cierre."

\end{lyrics}

\section{Resto qui}
\tag{thomas-entourage}\tag{sylvia-lai}\tag{italian-pop}

\subsection{Notes}
Based on the \emph{Thomas' Entourage} fictional universe by Carlos Thompson.\\
Character: Sylvia Lai.\\
Co-written between Carlos Thompson and Claude. High interaction.

\subsection{Style}
pop italiano, 2020s production\\
tempo: mid-uptempo, ♩≈ 102

\begin{lyrics}

\songpart{Verse 1}
\annotation{Mezzo-soprano vocals}\\
Barcellona la mattina odora di sale e caffè,\\
la luce entra diversa — più orizzontale, più mia.\\
Ho cercato questa sensazione in troppe città\\
e adesso so riconoscerla quando arriva.

\songpart{Verse 2}
Cinque anni di aeroporti, fusi orari, cabine,\\
Tokyo d'inverno, Sydney d'agosto, il Golfo a maggio.\\
Il mondo è piccolo quando lo percorri abbastanza —\\
piccolo e bellissimo e non è casa.

\songpart{Pre-Chorus}
\annotation{Whisper}\\
Non sono finita qui —\\
ho scelto qui.\\
C'è differenza,\\
e la conosco adesso.

\songpart{Chorus}
Resto qui dove il mare si vede dalla finestra,\\
resto qui con le persone che ho scelto io.\\
Ho girato abbastanza per sapere cosa cerco —\\
questa città, questa luce, questa vita che è mia.\\
Resto qui — non per mancanza di altrove,\\
ma perché altrove l'ho già visto tutto.

\songpart{Verse 3}
Vicky e Erika quando ci sono, il solito casino,\\
Gabi ha portato le tele e ha scelto questa città.\\
Ho amici che non passano per nessun altro —\\
sono miei, e questo è nuovo, e conta.

\songpart{Pre-Chorus}
\annotation{Whisper}\\
Sto finendo qualcosa\\
che ha solo il mio nome.\\
Non so ancora esattamente dove porta —\\
ma è la prima strada che ho costruito da sola.

\songpart{Chorus}
Resto qui dove il mare si vede dalla finestra,\\
resto qui con le persone che ho scelto io.\\
Ho girato abbastanza per sapere cosa cerco —\\
questa città, questa luce, questa vita che è mia.\\
Resto qui — non per mancanza di altrove,\\
ma perché altrove l'ho già visto tutto.

\songpart{Bridge}
\annotation{Spoken word}\\
Thomas era la strada — l'ho già detto a me stessa,\\
l'ho scritto, l'ho capito, l'ho messo via.\\
Sei anni di mondo che non avrei visto altrimenti,\\
un'educazione che nessuna scuola avrebbe dato.\\
Lo porto senza peso e senza cerimonia —\\
era quello che era, e mi ha portata qui.

\songpart{Verse 4}
Nello specchio i capelli sono teal da sei anni —\\
non è più una scelta, è semplicemente la mia faccia.\\
Carbonia esiste, la nonna conta ancora i giorni,\\
e io la amo da una distanza che è quella giusta.

\songpart{Pre-Chorus}
\annotation{Whisper}\\
Ho ventitré anni\\
e so già molte cose.\\
Non so ancora tutto —\\
ma voglio scoprirlo da qui.

\songpart{Chorus}
Resto qui dove il mare si vede dalla finestra,\\
resto qui con le persone che ho scelto io.\\
Ho girato abbastanza per sapere cosa cerco —\\
questa città, questa luce, questa vita che è mia.\\
Resto qui — non per mancanza di altrove,\\
ma perché altrove l'ho già visto tutto.

\songpart{Chorus (ripetizione finale — aperta, più alta)}
Resto qui e non mi scuso di niente,\\
resto qui con tutto quello che sono diventata.\\
Il peccato, la strada, le città, le persone —\\
tutto quello che è successo in mezzo non si riassume.
\annotation{Spoken word}\\
Resto qui — e per la prima volta\\
qui significa qualcosa che ho scelto io.

\end{lyrics}

\section{Choosing Here}
\tag{thomas-entourage}\tag{sylvia-lai}\tag{italian-pop}

\subsection{Notes}
Based on the \emph{Thomas' Entourage} fictional universe by Carlos Thompson.\\
Character: Sylvia Lai.\\
Interactive translation by Claude of \emph{Resto qui}.

\subsection{Style}
pop italiano, 2020s production\\
tempo: mid-uptempo, ♩≈ 102

\begin{lyrics}

\songpart{Verse 1}
\annotation{Female vocals}
Barcelona in the morning smells of salt and coffee,\\
the light falls differently—wider, warmer, more hers.\\
She chased that feeling through too many cities,\\
and now she knows it the moment it arrives.

\songpart{Verse 2}
Five years of airports, time zones, boarding gates,\\
Tokyo in winter, Sydney in August, the Gulf in May.\\
The world grows small when you have crossed it enough—\\
small, and so beautiful, but never home.

\songpart{Pre-Chorus}
\annotation{Whisper}
She did not end up here—\\
she chose this place.\\
There is a difference,\\
and she knows it now.

\songpart{Chorus}
Sylvia stays where the sea fills up her window,\\
she stays with people chosen by her heart.\\
She’s traveled far enough to know her longing—\\
this town, this light, this life that is her own.\\
She stays—not for a lack of other places,\\
but since she’s already beheld them all.

\songpart{Verse 3}
Vicky and Erika bring the usual chaos,\\
Gabi brought her canvases and chose this town.\\
She has friends who belong to no one else—\\
they’re hers, and that is new, and it matters.

\songpart{Pre-Chorus}
\annotation{Whisper}
She’s finishing something\\
that carries only her name.\\
She doesn’t quite know where it leads—\\
but it’s the first path she built with her own hands.

\songpart{Chorus}
Sylvia stays where the sea fills up her window,\\
she stays with people chosen by her heart.\\
She’s traveled far enough to know her longing—\\
this town, this light, this life that is her own.\\
She stays—not for a lack of other places,\\
but since she’s already beheld them all.

\songpart{Bridge}
\annotation{Spoken word}
Thomas was the road—she’s told herself already,\\
wrote it down, understood it, laid it to rest.\\
Six years of a world she’d have never seen otherwise,\\
an education no school could ever give.\\
She carries it lightly, no drama, no weight—\\
it was what it was, and it brought her here.

\songpart{Verse 4}
Six years of teal hair in the mirror now—\\
no longer a choice, it’s simply her face.\\
Carbonia exists, her grandmother counts the days,\\
and Sylvia loves her from a distance that is just right.

\songpart{Pre-Chorus}
\annotation{Whisper}
She is twenty-three\\
and already knows so much.\\
She doesn’t know everything yet—\\
but she wants to discover it all from here.

\songpart{Chorus}
Sylvia stays where the sea fills up her window,\\
she stays with people chosen by her heart.\\
She’s traveled far enough to know her longing—\\
this town, this light, this life that is her own.\\
She stays—not for a lack of other places,\\
but since she’s already beheld them all.

\songpart{Chorus (ripetizione finale — aperta, più alta)}
Sylvia stays here and makes no apologies,\\
she stays here with all she has become.\\
The mistakes, the roads, the cities, the people—\\
what happened between them cannot be summed up.
\annotation{Spoken word}
She stays here—and for the very first time\\
here means a place she chose for herself.

\end{lyrics}

\section{Melgar Table}
\tag{thomas-entourage}\tag{gabriele-wiater}\tag{blues}

\subsection{Notes}
Original work.
Narcocorrido narrative exploring themes of migration and economic desperation.

\subsection{Style}
Blues rock with a mid-tempo shuffle and loose backbeat, warm guitar bite and organ swells; verse rides a sparse groove with brushed snare, bass, and one overdriven riff, pre-chorus opens with tambourine and rising fills, chorus hits wider with gang vocals and stop-time turns. Lead vocal is gritty and close-mic with doubled lines, cracked ad-libs, and short delay throws on the hook. Ear candy: porch-noise fade-in, handclaps into the chorus, and a final guitar squeal lift. Mix is warm, roomy, and lived-in., blues, slow, rock, simple

\begin{lyrics}
\songpart{Verse 1}
My sister came in first\\
Dust on her shoes\\
My folks rolled up\\
From Barranquilla too

Tired in the eyes\\
Still laughing that way\\
Set the bags down slow\\
And the house felt okay

\songpart{Pre-Chorus}
No big show\\
No grand plan\\
Just the table\\
And every hand

\songpart{Chorus}
Easy New Year's Eve\\
Easy, easy\\
Pass that plate\\
And pass that breeze

Easy New Year's Eve\\
Easy, easy\\
We made a home\\
Out of what we need

\songpart{Verse 2}
Then you came down\\
From Bogotá\\
With the kids beside you\\
And Pilar in the car

Gabi flew in\\
Sylvia and me\\
Melgar by night\\
Was where we all should be

\songpart{Pre-Chorus}
Rooms kept changing\\
Doors swung wide\\
One more blanket\\
One more side

\songpart{Chorus}
Easy New Year's Eve\\
Easy, easy\\
Pass that plate\\
And pass that breeze

Easy New Year's Eve\\
Easy, easy\\
We made a home\\
Out of what we need

\songpart{Bridge}
Sylvia by my sister\\
Till my aunt arrived\\
Then the beds shifted\\
Like the rooms knew life

Pilar pitched her tent\\
Right out in the yard\\
Claiming her corner\\
Soft and hard

\songpart{Chorus}
Easy New Year's Eve\\
Easy, easy\\
Chicken, pork skin\\
Moussaka heat

Easy New Year's Eve\\
Easy, easy\\
Shared work, shared night\\
That was all we need

\songpart{Outro}
No service\\
No outside hand\\
Just family motion\\
And a crooked plan

I took slow walks\\
With Gabi near\\
Then slipped in the pool\\
And let it all clear
\end{lyrics}

\section{The Pirate's Ledger}
\tag{original-work}\tag{musical-theater}

\subsection{Notes}
Original work.
Musical theater-style narrative song drawn from historical maritime themes.

\subsection{Style}
Musical theater showtune with folk-pop influences. Baritone male vocals deliver a narrative performance with clear diction and theatrical phrasing. The arrangement features a bright acoustic guitar strumming eighth-note patterns, a melodic accordion providing counter-melodies and sustained chords, and a double bass playing on beats 1 and 3. A light percussion section includes a tambourine on the backbeat and subtle shaker. The tempo is 115 BPM in 4/4 time, likely in the key of G major. The harmonic structure follows a diatonic progression typical of contemporary musical theater, with a clean, dry production style that emphasizes vocal clarity and acoustic instrument separation.

\begin{lyrics}
\songpart{Verse 1}
\annotation{strummed acoustic guitar, double bass, accordion}\\
They called my father a pirate king\\
But piracy is a point of view\\
We farmed the hills in the northern spring\\
And traded goods when the season's blue \annotation{tambourine enters on backbeat}\\
Ransom was a form of trade\\
A merchant's contract signed in fear\\
We took the gold, the captive stayed\\
And everyone went home from here
\end{lyrics}

\section{Guided Hands}
\tag{original-work}\tag{musical-theater}

\subsection{Notes}
\songpart{Verse 1}
My phone screams red again\\
I swipe it blind\\
Three alarms, four alarms\\
Still I cut the line

You call to check my steps\\
I tell you I'm on track\\
I say "yeah, I'm getting there"\\
Then I never call you back

\songpart{Pre-Chorus}
I can look busy\\
I can look clean\\
Folders in a row\\
Like a perfect machine

But I fold in half\\
When nobody sees\\
Hide the delay\\
Wear the fake receipts

\songpart{Chorus}
I need guided hands\\
I need guided hands\\
I need guided hands\\
When I disappear (guided hands)

I need a room with rules\\
A face that stays\\
A place that knows my lies\\
And keeps me in the frame

\songpart{Verse 2}
I sit in bright white rooms\\
With a clipboard stare\\
I move the paper around\\
Like I meant to care

Tabs in a neat little line\\
Cursor parked and still\\
I fake the motion well\\
But I chase the will

I build little walls of proof\\
Just to feel okay\\
Then slip out through the crack\\
And waste another day

\songpart{Pre-Chorus}
I don't need a slogan\\
I don't need a speech\\
I need a held-closed door\\
And someone in reach

A group with names like mine\\
A guide who won't drift\\
Something steady enough\\
To catch my hands when I split

\songpart{Chorus}
I need guided hands\\
I need guided hands\\
I need guided hands\\
When I disappear (guided hands)

I need a room with rules\\
A face that stays\\
A place that knows my lies\\
And keeps me in the frame

\songpart{Bridge}
Maybe I'm not broken\\
Maybe I'm just lost\\
In a maze of yes and no\\
And a count of what it cost

So let me sit in truth\\
Let me say it plain\\
I keep building alibis\\
Then I wear their chains

\songpart{Final Chorus}
I need guided hands\\
I need guided hands\\
I need guided hands\\
When I disappear (guided hands)

I need a room with rules\\
A face that stays\\
A place that knows my lies\\
And keeps me in the frame

I need guided hands\\
I need guided hands\\
I need guided hands\\
And some way to stay (guided hands)

\subsection{Style}
Musical theater showtune with folk-pop influences. Baritone male vocals deliver a narrative performance with clear diction and theatrical phrasing. The arrangement features a bright acoustic guitar strumming eighth-note patterns, a melodic accordion providing counter-melodies and sustained chords, and a double bass playing on beats 1 and 3. A light percussion section includes a tambourine on the backbeat and subtle shaker. The tempo is 115 BPM in 4/4 time, likely in the key of G major. The harmonic structure follows a diatonic progression typical of contemporary musical theater, with a clean, dry production style that emphasizes vocal clarity and acoustic instrument separation.

\begin{lyrics}
\end{lyrics}

\section{Last Night's Extranger}
\tag{thomas-entourage}\tag{natalia-kyou-park}\tag{dubstep}

\subsection{Notes}
Based on the \emph{Thomas' Entourage} fictional universe by Carlos Thompson.\\
Character: Natalia Kyou Park.\\
Co-written between Carlos Thompson and Claude. High interaction.

\subsection{Style}
Dubstep with a half-time wobble groove, syncopated drums, growling bass drops, and crisp snare snaps; soft female lead; verse rides sparse sub pulses and clipped synth plucks, pre-chorus strips to tension risers and vocal chops, chorus hits with serrated bass hooks and stacked gang chants. Lead vocal stays close-mic and smoky in verses with doubled key lines, then wider shouted doubles on the hook. Ear candy: reversed swells, metallic ticks, and a filtered build before each drop. Bright, punchy, and high-gloss., abstract, dubstep, stylish, pixie, light, piercing, warm, theme, bright, sophisticated

\begin{lyrics}

\songpart{Verse 1}
New city, new badge,\\
new floor I don't know yet.\\
Milan in July,\\
still light at eight PM.

Suitcase half-open,\\
mirror says I'm ready.\\
Biggest job I've had\\
and no one here to tell.

\songpart{Pre-Chorus}
\annotation{airy, self-satisfied}\\
So I went downstairs\\
just to feel the air.\\
One drink, maybe two,\\
I wasn't going anywhere.

\songpart{Verse 2}
He sat down easy,\\
didn't ask my title.\\
Talked about the city\\
like he owned a piece of it.

No last names,\\
no LinkedIn,\\
just a good night\\
that got longer than I planned.

\songpart{Pre-Chorus}
So I went downstairs\\
just to feel the air.\\
One drink, maybe two,\\
I wasn't going anywhere.

\songpart{Chorus}
\annotation{cold drop, dry delivery}\\
Walked into the room,\\
nine AM, first day.\\
He's at the head of the table.\\
I'm on his payroll.

Last night's stranger\\
signs my checks now.\\
Milan, you absolute—\\
I'm on his payroll.

\songpart{Verse 3}
Nobody saw it.\\
I kept my face still.\\
Opened my notebook,\\
clicked my pen twice.

Is this funny?\\
Yeah, this is funny.\\
Or it will be\\
in about a year.

\songpart{Pre-Chorus}
\annotation{loaded, slower}\\
So I'd gone downstairs\\
just to feel the air.\\
One drink. Maybe two.\\
I had no idea.

\songpart{Chorus}
Walked into the room,\\
nine AM, first day.\\
He's at the head of the table.\\
I'm on his payroll.

Last night's stranger\\
signs my checks now.\\
Milan, you absolute—\\
I'm on his payroll.

\songpart{Bridge}
\annotation{rapped, close-mic, dry, fast}\\
Michael's in Los Feliz,\\
watching the same shows,\\
texting me goodnight\\
like he always does.\\
We don't talk about\\
what happens when we're gone.\\
We don't talk about it\\
because we don't have to.\\
He's got his work trips,\\
I've got mine.\\
Six thousand miles\\
is its own kind of time.\\
I love him, yeah,\\
that's not the question.\\
The question is\\
what just walked into my morning.

\songpart{Verse 4}
He didn't blink either.\\
Passed me the agenda.\\
Professional, clean,\\
like I imagined it.

Same kind of person.\\
Same kind of distance.\\
This might actually be\\
the least complicated thing.

\songpart{Pre-Chorus}
\annotation{quiet resolve}\\
So I went back downstairs\\
every night that week.\\
One drink, maybe two.\\
I knew exactly where I was going.

\songpart{Chorus}
Walked into the room,\\
nine AM, first day.\\
He's at the head of the table.\\
I'm on his payroll.

Last night's stranger\\
signs my checks now.\\
Milan, you absolute—\\
I'm on his payroll.

\songpart{Outro}
\annotation{synth fade, half-tempo}\\
I'm on his payroll.\\
I'm on his payroll.

\end{lyrics}

\section{Kyou Park}
\tag{thomas-entourage}\tag{natalia-kyou-park}\tag{electronic-pop}

\subsection{Notes}
Based on the \emph{Thomas' Entourage} fictional universe by Carlos Thompson.\\
Character: Natalia Kyou Park.\\
Co-written between Carlos Thompson and Claude. High interaction.

\subsection{Style}
female vocals, third person portrait, electronic pop, cold dossier verses,\\
warm intimate pre-choruses,\\
sparse synth verses, close-mic bridge, steady 4-beat grid, 105 BPM

\begin{lyrics}

\songpart{Verse 1}
Green eyes, folded lid,\\
three bloodlines, no flag.\\
Brown hair, green locks,\\
the choker matches.

Black shirt, pendant earrings,\\
nose ring, clean sneakers.\\
Any city fits her\\
like she always lived there.

\songpart{Pre-Chorus}
\annotation{warm, interior}\\
She knows the difference\\
between passing through\\
and being somewhere.\\
She learned it by doing both.

\songpart{Chorus}
Reads the system,\\
finds the line.\\
Makes it look easy,\\
every time.

Aerospace. Coast.\\
Head of the table.\\
Natalia didn't explain it.\\
She just arrived.

\songpart{Verse 2}
First stop: aerospace.\\
Engineers, tight specs.\\
She made it legible.\\
They trusted the work.

Then a coastline project,\\
chief of communications.\\
Two years on the ground.\\
Then she ran it.

\songpart{Pre-Chorus}
\annotation{quiet authority}\\
She doesn't raise her voice\\
to hold the room.\\
She already decided\\
before anyone sat down.

\songpart{Chorus}
Reads the system,\\
finds the line.\\
Makes it look easy,\\
every time.

Aerospace. Coast.\\
Head of the table.\\
Natalia didn't explain it.\\
She just arrived.

\songpart{Verse 3}
Michael's in LA,\\
same house, same time.\\
Technically exclusive —\\
that covers the defined life.

The traveling one runs parallel.\\
No guilt examined.\\
She calls it recalibration.\\
He probably does too.

\songpart{Pre-Chorus}
\annotation{intimate, matter-of-fact}\\
It's not betrayal\\
if it doesn't touch the thing.\\
She knows where home is.\\
She just doesn't always sleep there.

\songpart{Chorus}
Reads the system,\\
finds the line.\\
Makes it look easy,\\
every time.

Aerospace. Coast.\\
Head of the table.\\
Natalia didn't explain it.\\
She just arrived.

\songpart{Bridge}
\annotation{sparse synth, close-mic, slower}\\
The green was chosen.\\
The choker was chosen.\\
The face just arrived\\
from three places at once.

Serious at work,\\
easy everywhere else.\\
Same person, both modes.\\
Both genuine.

She doesn't sentimentalize.\\
She doesn't perform.\\
Complex systems, made legible.\\
Always has.

\songpart{Chorus}
\annotation{final, slightly broader}\\
Reads the system,\\
finds the line.\\
Makes it look easy,\\
every time.

Aerospace. Coast.\\
Head of the table.\\
Natalia didn't explain it.\\
She just arrived. 
(She just arrived)

\end{lyrics}

\section{Corona de alambre}
\tag{ricochets}\tag{nestor-mojica}\tag{narcocorrido}

\subsection{Notes}
Based on characters from \emph{Ricochets} by Carlos Thompson.\\
Character: Néstor Mojica, "Rey Pilos."\\
Co-written by Carlos Thompson, Claude, and Gemini. Highly guided and interacted.

\subsection{Style}
Narcocorrido norteño with brisk bajo sexto, tuba, accordion fills and corridos tumbados punch; verse rides tense but steady with tight rhyming lines, pre-chorus tightens on snare rolls and shouted asides, chorus lands in a chantable minor-key hook. Lead vocal is rough and nasal with doubled ends, group shouts on the hook word, short accordion turns between phrases, and a dusty, bright mix with dry kick, biting brass, and a final lift into an ominous break.

\begin{lyrics}

\songpart{Spoken Intro}
\annotation{Music fades in with a heavy, somber guitar or dark synth bass}\\
Esta es la historia de Rey Pilos, ganadero, narco y paramilitar.\\
El hombre que se quiso legalizar y conservar la operación."

\songpart{Verse 1}
Mojica era Rey Pilos,\\
con la mira en salir.\\
Quiso entrar por la rendija\\
que el gobierno abrió al rendir.

Se sentó con Don Pablo,\\
para mover el tablero.\\
Legalizar la historia\\
y cambiarse el sombrero.

\songpart{Pre-Chorus 1}
Pero antes de dar el paso,\\
había cuentas que limpiar.\\
Los Diablos de Don Pablo\\
ya buscaban dónde golpear.

La calle pidió silencio,\\
y el negocio, precisión.\\
Quien estorba en la jugada\\
sale del mapa, sin razón.

\songpart{Chorus}
Reino de papel, corona de alambre,\\
se compra el futuro matando el hambre.\\
Rey Pilos se va con un pacto en la mano,\\
pero la ley no olvida la sangre del hermano.\\
Nadie se salva, la cuenta es exacta:\\
cuando el pasado cobra, la tumba es la que pacta.

\songpart{Verse 2}
En un parking de Miami\\
Samuel Gunther cayó;\\
entre ráfagas y balas\\
madre e hija vieron caer.\\
Héctor Mendoza llegaba,\\
la gerencia de Inter-Flota asumió,\\
pero la ruta del sureste\\
ya no se pudo defender.

Mojica quemó ese puente,\\
lo soltó sin compasión.\\
Si la ruta ya era veneno,\\
mejor tumbarla al montón.

\songpart{Pre-Chorus 2}
\annotation{Conversational}\\
Lejos de Miami, la guerra estalló,\\
y a los socios de Mojica la muerte alcanzó:\\
un camión bomba a Rodríguez le frenó el carruaje,\\
y en Colombia a Arbeláez le cobraron el peaje.

Morocho movía el plano, Don Pablo dio la señal,\\
usando al agente Bard como ficha electoral.\\
Le entregaron las sobras a la ley federal,\\
limpiando el terreno en un juego político y criminal.

\songpart{Chorus}
Reino de papel, corona de alambre,\\
se compra el futuro matando el hambre.\\
Rey Pilos se va con un pacto en la mano,\\
pero la ley no olvida la sangre del hermano.\\
Nadie se salva, la cuenta es exacta:\\
cuando el pasado cobra, la tumba es la que pacta.

\songpart{Verse 3}
Lo que no vio Rey Pilos\\
fue el archivo en la oficina:\\
el fundador de Inter-Flota\\
ya sabía la cocina.

Y un abogado de auditoría\\
halló la red del sureste.\\
Cada cobro, cada ruta,\\
cada puerta y cada oeste.

No querían quemar el tronco,\\
solo las ramas del jardín,\\
pero el informe completo\\
puso a la DEA en el festín.

\songpart{Bridge}
Se firmó con traje serio,\\
pero olía a confesión.\\
La legalidad de un reino\\
a veces cuesta la prisión.

Y en la mesa de los nombres\\
nadie compra el corazón.\\
Porque el orden que se inventa\\
también paga su traición.

\songpart{Chorus}
Reino de papel, corona de alambre,\\
se compra el futuro matando el hambre.\\
Rey Pilos se va con un pacto en la mano,\\
pero la ley no olvida la sangre del hermano.\\
Nadie se salva, la cuenta es exacta:\\
cuando el pasado cobra, la tumba es la que pacta.

\songpart{Outro, (The music slows down, leaving just a heavy drum beat)}
Dos años dieron la vuelta... y llegó la extradición.\\
El gobierno los mandó lejos... con la ley en el avión.
\annotation{Fade out}

\end{lyrics}

\section{¿Aceptás el trato?}
\tag{ricochets}\tag{fabian-muller}\tag{rock-en-espanol}

\subsection{Notes}
Based on characters from \emph{Ricochets} by Carlos Thompson.\\
Character: Fabián Ernesto Müller Char.\\
Interactively fixed between Carlos Thompson and Gemini, over an autogenerated lyrics by Suno, in turn based on a \emph{Ricochets' Companion Guide} story written by Carlos Thompson. (original Suno lyrics was too Mexican + Coliche, which didn't fit the character.)

\subsection{Style}
Rock en español with blues grit and punk energy, medium-fast driving groove; verse rides tight bass, clipped guitar and tom-heavy percussion, pre-chorus opens with chord stabs and keyboard swells, chorus hits with melodic backing harmonies and smooth tenor doubles. Bridge drops to a bluesy half-time shuffle then surges back. Close-mic velvet tenor vocal, soulful ad-libs, warm delay on the hook, gritty punchy instrument mix, smooth vocal mix, male vocals, clean melodic vocals, world, blues, punk, rock

\begin{lyrics}

\songpart{Verse 1}
Salí de la escuela\\
con ganas de hacer\\
un nombre limpio,\\
un imperio de bien.

Miraba los mapas,\\
logística fiel,\\
moviendo la carga,\\
creciendo la red.

Cartagena al hombro,\\
Miami en la piel,\\
rutas perfectas\\
que armé en el papel.

Soñaba una empresa\\
más grande también:\\
poner a Colombia\\
a marchar al nivel.

\songpart{Pre-Chorus}
Pero el negocio\\
empezó a torcer,\\
manos extrañas\\
firmando el poder.

Cuentas fantasmas\\
difíciles de ver,\\
rutas fantasmas\\
difíciles de creer.

Contenedores\\
que abrieron la red,\\
un virus callado\\
que entró sin querer.

\songpart{Chorus}
¿Aceptás el trato?\\
¿Mirás para allá?\\
No jugués al limpio\\
si el juego cambió.

¿Aceptás la cuota?\\
¿O vas a hablar?\\
Aquí manda el jefe,\\
aquí no hay plan B.

\songpart{Verse 2}
Los libros decían\\
que todo iba bien,\\
ingresos perfectos,\\
las leyes de ley.

Yo quise cegarme,\\
no quise entender,\\
si abría los ojos\\
podía caer.

Los socios de siempre\\
me hicieron dudar,\\
comprando silencios\\
con tanto poder.

Y yo en mi soberbia,\\
queriendo creer:\\
"Si yo no firmé,\\
nada tengo que ver".

\songpart{Pre-Chorus}
Hasta que una noche\\
cayó la pared,\\
la DEA en la puerta,\\
sin nada que ver.

Doscientos kilos\\
que el barco cargó,\\
me dejaron solo,\\
frente a la ley.

\songpart{Chorus}
¿Aceptás el trato?\\
¿Mirás para allá?\\
No jugués al limpio\\
si el juego cambió.

¿Aceptás la cuota?\\
¿O vas a hablar?\\
Aquí manda el jefe,\\
aquí no hay plan B.

\songpart{Bridge}
El abogado fue claro,\\
sin darme color:\\
"Tu firma está puesta\\
en cada estibador".

"Si entregás los nombres,\\
podés negociar,\\
pero si callás,\\
te van a enterrar".

Sentí que la vida\\
perdió su valor,\\
atrapado en medio\\
del miedo y el honor.

Mi sueño educado,\\
mi gran ambición,\\
muriendo en la celda\\
de mi propia omisión.

\songpart{Final Chorus}
¿Aceptás el trato?\\
¿Mirás para allá?\\
Se acabó el engaño,\\
no hay cómo fingir.

¿Aceptás la culpa?\\
¿O vas a hablar?\\
O hundís tu silencio,\\
o aprendés a vivir.

¿Aceptás el trato?\\
¿O te vas a quebrar?\\
Si no sos de ellos,\\
te van a borrar.

\end{lyrics}

\section{Vaina que crece}
\tag{original-work}\tag{art-pop}\tag{spanish-lyrics}

\subsection{Notes}
Original work.
Theme: Mental health, accountability, and the need for structured support systems.

\subsection{Style}
Art pop electrónico con pulso medio, groove sincopado y un crecimiento gradual de tensión; verso en capas mínimas con bajo filtrado y piano seco, pre-coro abre espacio con armonías aireadas y percusión más marcada, coro entra ancho con sintes brillantes y baterías firmes. Voz íntima al frente con dobles suaves, coros apilados en el gancho y pequeñas colas de delay en palabras clave. Riser reverso antes del coro, chasquidos digitales entre frases y un estallido de textura en el final. Mezcla limpia, grande y luminosa.

\begin{lyrics}
\songpart{Verse 1}
Una pequeña vaina comienza\\
y poco a poco adquiere forma\\
se elabora\\
en silencio

Llega la distancia\\
el cambio\\
el recuerdo\\
de un obscuro viaje al Despertar

La inspiración que cede al Græville\\
y se pierde\\
en un par de experimentos estúpidos\\
hasta que duele más la espera

\songpart{Pre-Chorus}
Y vuelve a moverse\\
bajo la piel\\
vuelve a romperse\\
vuelve a nacer

La crisis llama\\
sin pedir perdón\\
y lo que era polvo\\
toma voz

\songpart{Chorus}
Una nueva historia\\
un nuevo sueño\\
una nueva historia\\
crece en el pecho

Que crece\\
que engrandece\\
que rompe el miedo\\
y vuelve a empezar

Una nueva historia\\
un nuevo sueño\\
una nueva historia\\
viene a estallar

\songpart{Verse 2}
Perdida en un abismo electrónico\\
la tarde se queda sin nombre\\
y todo gira\\
sin dirección

Entre chispas\\
y gestos torpes\\
se arma el vacío\\
se parte el sol

Y de repente\\
sin pedir permiso\\
la idea respira\\
se levanta sola

\songpart{Pre-Chorus}
Se junta el pulso\\
con el dolor\\
se enciende el borde\\
del corazón

La crisis llama\\
sin pedir perdón\\
y lo que era ruina\\
toma voz

\songpart{Chorus}
Una nueva historia\\
un nuevo sueño\\
una nueva historia\\
crece en el pecho

Que crece\\
que engrandece\\
que rompe el miedo\\
y vuelve a empezar

Una nueva historia\\
un nuevo sueño\\
una nueva historia\\
viene a estallar

\songpart{Bridge}
Y llega la distancia\\
ya no pesa igual\\
el recuerdo no manda\\
puede soltar

De un oscuro viaje al Despertar\\
nace la ruta\\
nace el umbral

\songpart{Final Chorus}
Una nueva historia\\
un nuevo sueño\\
una nueva historia\\
crece en el pecho

Que crece\\
que engrandece\\
que rompe el miedo\\
y vuelve a empezar

Una nueva historia\\
un nuevo sueño\\
una nueva historia\\
viene a estallar
\end{lyrics}

\section{Monoestable}
\tag{original-work}\tag{alternative-rock}\tag{spanish-lyrics}

\subsection{Notes}
Original work in Spanish.
Explores state dynamics and internal oscillation; themed around instability and emotional transformation.

\subsection{Style}
Dark alternative pop-rock in Spanish with a tense mid-tempo pulse, jagged palm-muted guitars, punchy toms, and a low synth drone; verse stays tight and claustrophobic, pre-chorus opens with rising backing vocals and filtered percussion, chorus lands with a stark shouted hook and layered doubles, bridge strips to voice and bass before the final lift. Close-mic lead vocal, dry and intimate in verses, wider harmonies and short delay throws on the hook word, with vinyl crackle and reverse swells between sections; bright top end over a heavy, compressed low end.

\begin{lyrics}
\songpart{Verse 1}
Teoría del chaos, teoría de catástrofes\\
es el gatillo que dispara al monoestable\\
¿Cómo se puede estar bien?\\
¿Cómo se puede caer así?

Sombras que se cierran\\
sobre la mente\\
que encierran, que aprisionan\\
que revientan

\songpart{Pre-Chorus}
Y no se abre la puerta\\
y no baja la tensión\\
todo gira hacia adentro\\
todo rompe alrededor

\songpart{Chorus}
Y es ella\\
siempre ella\\
Ella y el vacío\\
Ella y el vacío

Y es ella\\
siempre ella\\
Me deja sin nombre\\
me deja sin hilo

\songpart{Verse 2}
Lleva el pulso en la nuca\\
como un vidrio bajo piel\\
me promete que regreso\\
y me pierde otra vez

En la esquina de mis días\\
se me quiebra la razón\\
y lo que callo me muerde\\
desde el fondo del montón

\songpart{Pre-Chorus}
Y no se abre la puerta\\
y no baja la tensión\\
todo gira hacia adentro\\
todo rompe alrededor

\songpart{Chorus}
Y es ella\\
siempre ella\\
Ella y el vacío\\
Ella y el vacío

Y es ella\\
siempre ella\\
Me deja sin nombre\\
me deja sin hilo

\songpart{Bridge}
Dime dónde guardo el miedo\\
si me ocupa todo el cuerpo\\
dime cómo sigue vivo\\
lo que nace de un incendio

\songpart{Chorus}
Y es ella\\
siempre ella\\
Ella y el vacío\\
Ella y el vacío

Y es ella\\
siempre ella\\
Me deja sin nombre\\
me deja sin hilo
\end{lyrics}

\section{E Numa Lah}
\tag{original-work}\tag{experimental-pop}

\subsection{Notes}
Original work.
Experimental composition with stylized, non-semantic vocals and abstract sound design.

\subsection{Style}
Experimental alt-pop with syncopated handclaps and a rubbery bass pulse; verse stays sparse with dry percussion and a single looping synth motif, pre-chorus lifts with filtered vocal stacks and a rising tom pattern, chorus lands wider with doubled lead vocals and a chanty call-and-response on the hook. Add ear candy with reversed swells, quick tape-stop moments, and sparkling chimes between phrases. Bright, tight mix with a playful, slightly surreal edge.

\begin{lyrics}
\songpart{Verse 1}
E dem chi laim siman de yacken\\
e vis se lah direch vi, direch vi\\
numa lah e numa lah...\\
e numa lah!

\songpart{Pre-Chorus}
Laim de inkera bah\\
se in ox mint de baiya

\songpart{Chorus}
E numa lah\\
E numa lah\\
Limpter vi, limpter vii\\
e numa lah

E numa lah\\
E numa lah\\
Limpter vi, limpter vii\\
e numa lah

\songpart{Verse 2}
E dem chi laim siman de yacken\\
e vis se lah direch vi, direch vi\\
numa lah e numa lah...\\
e numa lah!

\songpart{Chorus}
E numa lah\\
E numa lah\\
Limpter vi, limpter vii\\
e numa lah

E numa lah\\
E numa lah\\
Limpter vi, limpter vii\\
e numa lah
\end{lyrics}

\section{Mirá}
\tag*{1996}\tag{original-work}\tag{musical-theater}

\subsection{Notes}
Original work from 1996.
Musical theater piece using conversational narrative framing.

\subsection{Style}
Musical theater showtune with folk-pop influences. Baritone male vocals deliver a narrative performance with clear diction and theatrical phrasing. The arrangement features a bright acoustic guitar strumming eighth-note patterns, a melodic accordion providing counter-melodies and sustained chords, and a double bass playing on beats 1 and 3. A light percussion section includes a tambourine on the backbeat and subtle shaker. The tempo is 115 BPM in 4/4 time, likely in the key of G major. The harmonic structure follows a diatonic progression typical of contemporary musical theater, with a clean, dry production style that emphasizes vocal clarity and acoustic instrument separation.

\begin{lyrics}

\songpart{Verse 1}
\annotation{strummed acoustic guitar, double bass, accordion}
Amiga, ¿puedo llamarte amiga? mirá! tengo un problema, tú lo conoces o lo has oído en parte, estoy seguro.  Resulta que, hace tiempo ya, encontré a una persona, tú la conoces, su nombre; si no lo sabes lo descubrirás. \annotation{tambourine enters on backbeat} Hace algún otro tiempo descubrí que la amaba, o al menos eso creí, o, ¿sabes?, creo que ni eso creí pero lo hice.  Y también hace algún tiempo, se lo hice saber.  No sé si me creyó 

\end{lyrics}

\section{Si pienso}
\tag{original-work}\tag{trance}\tag{spanish-lyrics}

\subsection{Notes}
Original work by Carlos Thompson.
In Spanish. Trance meditation on memory and obsessive thought patterns.

\subsection{Style}
trance, latin pop, soprano lead vocals, latin-american pronunciation, bassline melody lead, airy vocal phrasing, delay throws, harmonized chorus, sidechain pumping, glossy synth pads, arpeggiated trance plucks, four-on-the-floor kick, gated reverb snare, 128 BPM, bittersweet longing, neon melancholy

\begin{lyrics}
Si pienso en una persona; pienso en tu rostro...\\
Si pienso en una voz; busco la tuya en mi memoria...\\
Si pienso en una mirada; la imagino de tus ojos...\\
Si pienso en lo que nunca fue; tu imagen aparece, tus ojos, tu cabello, tu boca, tu piel...\\
Si pienso en qué deseo; primero estás tú, porque tú eres lo que más deseo...\\
Si pienso en mi dolor; recuerdo aquellas palabras que tú me dijiste, aquellos saludos que no respondiste...\\
Si pienso en mi vida; necesariamente tengo que pensar en ti, en las esperanzas que aún no he perdido, en las imágenes que de ti aún guardo, en unos ojos abiertos mirando a todas partes exceptuándome a mi, en una voz dulce con palabras ásperas, en un nombre que temo pronunciar, un rostro que no puedo olvidar, tantas cosas que desearía no recordar...\\
Pero no puedo...\\
Si pienso en cualquier cosa, ahí, visible o escondida en algún lado, ahí estás tú.
\end{lyrics}

\section{Iros de risk}
\tag{original-work}\tag{trance}\tag{spanish-lyrics}

\subsection{Notes}
Original work.
Experimental trance composition with fragmented narrative structure.

\subsection{Style}
trance, latin pop, soprano lead vocals, latin-american pronunciation, bassline melody lead, airy vocal phrasing, delay throws, harmonized chorus, sidechain pumping, glossy synth pads, arpeggiated trance plucks, four-on-the-floor kick, gated reverb snare, 128 BPM, bittersweet longing, neon melancholy

\begin{lyrics}
Iros de risk, nefa he sabido como entrar, solo vi y temí, e iros de risk.  Decidí y ahí se va.  Ce ça, aill be, je bi.\\
O.K., she's 'ere and aim 'ere.\\
Pero no será\\
no será en aim sure
'n aim' sure\\
sure\\
shit\\
shes 'ere\\
en' aim 'ere\\
ce nothink\\
nada!
\end{lyrics}

\section{Uneekwa En Urbo}
\tag{original-work}\tag{esperanto-lyrics}\tag{folk}

\subsection{Notes}
Original work.
Esperanto folk duet, a portrait of an urban figure presented through dual vocal perspectives.

\subsection{Style}
Esperanto folk duet with gentle 6/8 sway and fingerpicked acoustic guitar; verse 1 is intimate and close, verse 2 answers in a warmer lower harmony, pre-chorus opens with soft accordion and hand percussion, chorus blooms with stacked two-part harmonies and a repeated chantable name hook. Add occasional breathy oohs, a brief bridge with only guitar and low harmony, then final chorus with light tambourine, street-noise field textures, and a wide, warm mix, folk

\begin{lyrics}
\songpart{Verse 1}
\annotation{M}\\
En la centro ŝi iras,\\
Uneekwa laŭ la fluo.\\
Nigraj plektaĵoj falas,\\
super ŝultroj kiel nokto.

\annotation{F}\\
Nigra denim’ ĉe la kruroj,\\
kaj botoj ĝis la genuo.
Ĉiu paŝo diras klare:
ŝi portas sian propran lumon.

\songpart{Pre-Chorus}
\annotation{M}\\
Oliva bluzo glitas,\\
kvazaŭ venta silk’ tra tago.

\annotation{F}\\
Kaj ŝia rideto brilas,\\
tra la bruo de la placo.

\songpart{Chorus}
Uneekwa, Uneekwa
Ĉiuj turnas sin al ŝi.\\
Uneekwa, Uneekwa\\
Kun forto kaj ĉarmo ĉi.

\annotation{M}\\
Bruna rigardo, saĝa fajro,
\annotation{F}\\
ebonhaŭto, pura reĝ’,
(Bela Uneekwa)
\annotation{M}
Ŝi marŝas, ŝi ne petas
\annotation{F}\\
la urbo faras spacon tuj.

\songpart{Verse 2}
\annotation{F}\\
Sur la vango ludas kavetoj,\\
kaj la lipoj portas vinon.\\
Verda ombro ĉirkaŭ okuloj,\\
kvazaŭ ĝardeno post pluvo.

\annotation{M}\\
Naza ringo, brilaj orelringoj,\\
arĝentĉeno ĉe la klaviklo.\\
Maldekstre verda sako pendas,\\
dekstre lumas telefono.

\songpart{Pre-Chorus}
\annotation{F}
Ŝi nur rigardas ekranon,\\
kaj la mondo sekvas laŭ ŝi.

\annotation{M}\\
Inter paŝoj kaj fenestroj,
ŝia ĉeesto kantas pli.

\songpart{Chorus}
Uneekwa, Uneekwa
Ĉiuj turnas sin al ŝi.\\
Uneekwa, Uneekwa\\
Kun forto kaj ĉarmo ĉi.

\annotation{M}\\
Bruna rigardo, saĝa fajro,
\annotation{F}\\
ebonhaŭto, pura reĝ’,
(Bela Uneekwa)
\annotation{M}
Ŝi marŝas, ŝi ne petas
\annotation{F}\\
la urbo faras spacon tuj.

\songpart{Bridge}
\annotation{M}\\
Ni aŭdas ŝin en la trotuaro,\\
kvazaŭ tamburo sub la piedoj.

\annotation{F}\\
Ni sekvas ŝian trankvilon,\\
kaj la tago ŝanĝas tonon.

\songpart{Final Chorus}
Uneekwa, Uneekwa
Ĉiuj turnas sin al ŝi.\\
Uneekwa, Uneekwa\\
Kun forto kaj ĉarmo ĉi.

\annotation{M}\\
Bruna rigardo, saĝa fajro,
\annotation{F}\\
ebonhaŭto, pura reĝ’,
(Bela Uneekwa)
\annotation{M}
Ŝi marŝas, ŝi ne petas
\annotation{F}\\
la urbo faras spacon tuj.
\end{lyrics}

\section{Potsdam Leeres Büro}
\tag{thomas-entourage}\tag{gabriele-wiater}\tag{industrial-electronic}

\subsection{Notes}
From the \emph{Thomas' Entourage} universe.
Character: Gabriele Wiater, SVA student in New York.
Inner monologue and emotional landscape.

\subsection{Style}
industrial electronic, 165 BPM, female lead vocal, restrained delivery, gang chorus, overdriven power chords, palm-muted verse guitar, octave bass line, snare room slap, live kick bleed, bathroom ambience, tape saturation, spring reverb, mono drum image, clipped vocal delay, adrenalized defiance, lonely drift, driving eighth notes

\begin{lyrics}
SVA in NYC. Ich bin endlich in der letzten Phase meiner Reise—meines Traums.\\
SVA in NYC… warum fühlt es sich so einsam an?

Diese Stadt—chaotisch, feindselig, anonym.\\
Vier Jahre lang habe ich Sprache und Kommunikation studiert, ein Portfolio aufgebaut, und all das wirkt jetzt leer.

Dieses Gefühl, von Menschen umgeben zu sein—meine Leute an der NYU, meine neuen Leute an der SVA—und am Ende doch niemanden zu haben, auf den ich mich verlassen kann. Niemanden, der nach mir sieht.

Und Potsdam?\\
Die Niederlage wäre real, wenn ich jemals zurückkäme.

Aber bin ich mutig genug, durchzuhalten?

\end{lyrics}

\section{Sunlit Gem Music}
\tag{original-work}\tag{orchestral}\tag{epic}

\subsection{Notes}
Original work.
A meditation on romantic idealization drawn from a classical-style love poem.

\subsection{Style}
epic, 78 BPM, dramatic string ostinato, French horn chorales, low timpani rolls, concert tom accents, harp arpeggios, solo cello melody, cathedral hall reverb, tape saturation, wide stereo mix, halftime march pulse, suspended cymbal swells, bittersweet devotion, weary longing, orchestral crescendo, minor key, restrained brass swells

\begin{lyrics}
Please interpret this poem. What is it about?

My eyes, they feast upon your sunlit face,\\
A landscape where each feature seems a gem,\\
And in your voice, a music finds its place,\\
A symphony that lulls to diadem.\\
O dearest one, you hold the highest rank,\\
My chosen star, above the rest alight,\\
Yet even stars grow dim, their brilliance dank,\\
When viewed too oft by eyes that lose their sight.\\
For fickle is the heart, a restless thing,\\
It craves the novel, shuns the known too well,\\
And though love's flame with fervent heat may cling,\\
The embers die in ashes' ashen shell.\\
Forgive, my love, this flaw that stains my core,\\
I seek not solace elsewhere, but adore.

\end{lyrics}

\section{Soap Line Logic}
\tag*{2026}\tag{original-work}\tag{electronic-pop}

\subsection{Notes}
Original work from 2026.
Abstract electronic reflection on mundane ritual and meditative routine.

\subsection{Style}
Techno-trance with dubstep drops, four-on-the-floor pulse, wobble bass under a melodic female lead; verse stays lean and mechanical, pre-chorus filters up with clipped handclaps and rising synth shimmer, chorus opens wide with octave doubles and chanty backing vocals, bridge strips to a half-time sub pulse and breathy topline before the final drop. Glassy arps, water-drip percussion, reverse swells into each transition, bright punchy mix with heavy low-end and clean vocal polish., relaxing, techno, dubstep

\begin{lyrics}

\songpart{Verse 1}
I leave the plates\\
For later, for later\\
My mind says “scroll”\\
My hands say “maybe”\\
The sink looks full\\
Like a small, dull mountain\\
But nothing breaks\\
If I wait a minute

\songpart{Pre-Chorus}
Then I turn on the tap\\
And the room changes shape\\
Hot water, cold steam\\
I find my pace\\
One bowl, one fork\\
One little fix\\
I line it up\\
And I click

\songpart{Chorus}
Doing the dishes, doing the dishes\\
(doing the dishes)\\
Back in the rhythm, back in the rhythm\\
(back in the rhythm)\\
Soap line logic\\
Soap line logic\\
A small, steady peace\\
In the rinse and the drip

\songpart{Verse 2}
I watch the grease\\
Slip off the spoon\\
Keep the sponge fresh\\
Keep the knife clean too\\
Stack the plates\\
Like a game I can solve\\
Tall ones first\\
Then the cups in a row\\
It’s funny how\\
A simple routine\\
Turns down the noise\\
Inside of me

\songpart{Pre-Chorus}
Then I turn on the tap\\
And the room changes shape\\
Hot water, cold steam\\
I find my pace\\
One bowl, one fork\\
One little fix\\
I line it up\\
And I click

\songpart{Chorus}
Doing the dishes, doing the dishes\\
(doing the dishes)\\
Back in the rhythm, back in the rhythm\\
(back in the rhythm)\\
Soap line logic\\
Soap line logic\\
A small, steady peace\\
In the rinse and the drip

\songpart{Bridge}
Maybe that’s the point\\
Of every little chore\\
Not to thrill me\\
Just to pull me to the floor\\
What I avoid\\
Can hold me in place\\
Till the static in my head\\
Starts to fade

\songpart{Drop}
Doing the dishes\\
Doing the dishes\\
(doing the dishes)\\
Doing the dishes\\
Back in the rhythm\\
Back in the rhythm\\
(soap line logic)\\
A small, steady peace\\
In the rinse and the drip

\songpart{Chorus}
Doing the dishes, doing the dishes\\
(doing the dishes)\\
Back in the rhythm, back in the rhythm\\
(back in the rhythm)\\
Soap line logic\\
Soap line logic\\
A small, steady peace\\
In the rinse and the drip

\end{lyrics}

\section{Equals}
\tag{thomas-entourage}\tag{mariain-schumer}\tag{gabriele-wiater}\tag{pop-ballad}

\subsection{Notes}
From the \emph{Thomas' Entourage} universe.
Characters: Mariain Schumer and Gabriele Wiater.
Co-written with ChatGPT. High interaction.
Explores emotional recalibration and mutual growth in relationship transformation.

\subsection{Style}
pop ballad, electro-pop, female mezzo lead, close-mic confessional, Rhodes piano, fingerpicked acoustic guitar, fretless bass, brushed snare, analog synth pads, soft tom fills, tape saturation, plate reverb, 92 BPM, swung 8th groove, bittersweet release, late-night intimacy, lifted chorus, 90s adult contemporary

\begin{lyrics}

\songpart{Verse 1}
I met you broken, freshman year, paint still wet on your hands\\
Emotional wreck in a big city, lost in a foreign land\\
I told myself I’d be your guide, but part of me wanted more\\
Somewhere between savior and sinner, I opened up the door

\songpart{Verse 2}
You bloomed right in front of me, turned my project into pride\\
From late-night talks to tangled sheets, with you right by my side\\
We went and made it official, moved your things into my place\\
Too much forever, too much us, we started losing space

\songpart{Pre-Chorus}
Jealousy came crawling in, every name became a scar\\
What felt free when we were friends just hurt when we called it ours

\songpart{Chorus}
Now since Thomas walked in, I feel less suffocated by you\\
We’re finally real friends again, not pretending we’re brand new\\
You’re lighter, sharper, glowing like someone who met her match\\
I thought that I was your equal… but I’m learning how to let that go\\
Yeah, I’m learning how to let that go

\songpart{Verse 3}
You came back from Germany with new fire in your eyes\\
Told me ‘bout a stranger on a plane and Tallinn midnight skies\\
Wasn’t love at first, you said — it was someone who saw you whole\\
And damn, that look you gave him carved a new space in my soul

\songpart{Pre-Chorus}
Business you’re more focused, as a friend you cut so deep\\
As lovers you’re more honest, more alive than when you’re with me

\songpart{Chorus}
Now since Thomas walked in, I feel less suffocated by you\\
We’re finally real friends again, not pretending we’re brand new\\
You’re lighter, sharper, glowing like someone who met her match\\
I thought that I was your equal… but I’m learning how to let that go

\songpart{Bridge}
I remember Carroll Gardens, the soiree, summer heat\\
You watching him across the room, studying every beat\\
Asked me what I thought of him like my opinion still mattered\\
I smiled, said he seems dangerous… in all the ways that you’ve been shattered

\songpart{Verse 4}
We were never meant for cages, we were always meant to fly\\
I loved the version that I saved, but I love the one that said goodbye\\
Maybe love looks different now, maybe this is how it heals\\
Two old flames becoming real friends underneath the same sky

\songpart{Pre-Chorus}
Jealousy still whispers, but I’m learning how to hush it down\\
‘Cause seeing you this free feels like the truest thing we’ve found

\songpart{Chorus - Outro}
Since Thomas walked in… I feel less suffocated by you\\
We’re finally intimate friends, the way we always should’ve been\\
You shine brighter when you’re honest, when you’re standing in your truth\\
I thought that I was your equal…\\
But maybe love just needed room\\
Yeah, maybe love just needed room…\\
To let you find somebody new

\end{lyrics}

\section{East Rock / The VIllage}
\tag{thomas-entourage}\tag{mariain-schumer}\tag{pop-ballad}

\subsection{Notes}
From the \emph{Thomas' Entourage} universe.
Character: Mariain Schumer.
Co-written with ChatGPT. High interaction.
Portrait of a woman navigating two cities and life trajectories.

\subsection{Style}
pop ballad, soul inflections, R\&B phrasing, 105 BPM, female lead, close-mic verse, soft chest voice, melismatic pre-chorus, restrained belt chorus, electric piano voicings, muted Fender bass, brushed snare groove, fingerpicked nylon guitar, warm analog tape, plate reverb, stereo delay throws, string swells, bittersweet nostalgia, urban dusk

\begin{lyrics}

\songpart{Verse 1}
She grew up where the furnaces went cold,\\
read every shelf before she found the door,\\
and Yale received the hunger that she carried —\\
to name the thing before it changed once more.

\songpart{Verse 2}
She chose East Rock when others chose the city,\\
put down the kind of roots that feed on choice,\\
learned permanence is also an ambition,\\
and planted herself quietly in the noise.

\songpart{Pre-Chorus}
She took the train down every week,\\
could read a room by smell and tone,\\
moved through the openings like water —\\
the art world is a market, and she'd known.

\songpart{Chorus}
But there is quiet in East Rock at dusk,\\
the maples and the brownstones holding still,\\
and there is something in the Village musk —\\
\emph{the Village} — that no other city will.\\
Mariain keeps two doors, two keys, two lives,\\
the particular light, the particular street.

\songpart{Verse 3}
At first she took on artists like a game,\\
she knew the fault lines, knew the leverage well,\\
then Gabi came with different weight entirely —\\
she called it professional, and couldn't tell.

\songpart{Pre-Chorus}
She found her in Long Island City,\\
forgetting food, forgetting cold —\\
she went because the friendship held them,\\
and stayed because the mirror showed.

\songpart{Chorus}
But there is quiet in East Rock at dusk,\\
the maples and the brownstones holding still,\\
and there is something in the Village musk —\\
\emph{the Village} — that no other city will.\\
Mariain keeps two doors, two keys, two lives,\\
the particular light, the particular street.

\songpart{Bridge}
She learned to loosen what her gravity had cost,\\
watched Gabi follow her own sound —\\
but liberation finds its irony:\\
the freer Gabi grew, the less she could be found.

Then Thomas, easy charm, a steadying hand,\\
and Mariain watched with something like relief,\\
a jealousy too honest to deny,\\
too mild to call anything but brief.

\songpart{Verse 4}
The week has its reliable clean shape now,\\
the artists and the openings and the train,\\
East Rock receives her with its patient maples,\\
the Village gives her back her edge again.

\songpart{Pre-Chorus}
And then the visa question settled,\\
not contested, barely said —\\
Gabi chose the other gravity,\\
and took the warmth when she left.

\songpart{Chore - Outro}
There is quiet in East Rock at dusk,\\
the maples and the brownstones holding still,\\
and there is something in the Village musk —\\
\emph{the Village} — that no other city will.\\
But tonight the rooms are only rooms,\\
the keys turn on the ordinary air —\\
somewhere over water, Gabi blooms,\\
and the particular street leads only there.

\end{lyrics}

\section{Lilac-eyed Demon}
\tag{thomas-entourage}\tag{mariain-schumer}\tag{dubstep}

\subsection{Notes}
From the \emph{Thomas' Entourage} universe.
Character: Mariain Schumer.
Dubstep accusation and introspective reckoning with impact and culpability.

\subsection{Style}
dubstep, club, 140 BPM, female soprano lead vocals, vocoder lead, formant shifting, syncopated half-time groove, sub bass drops, clipped synth stabs, sidechain pumping, gated reverb snare, metallic percussion, sparse verse arrangement, stacked hook harmonies, retro-futurist sheen, nocturnal energy

\begin{lyrics}

Mariain — you've lived your life as a soulless (as a soulless)\\
Mariain — you crushed dreams and fed on them (fed on them)\\
Your lilac eyes were their perdition (their perdition)\\
Your flawless speech the way you got them (you got them) 

Mariain, lilac-eyed demon (lilac-eyed demon) 

Mariain — you've lived your life as a soulless (as a soulless)\\ 
Mariain — you went and fell in love with your prey (with your prey)\\
Then you couldn't handle the situation (the situation)\\
And your Gabrielè almost died there (almost died there)

Mariain, lilac-eyed demon (lilac-eyed demon) 

Mariain — you've lived your life as a soulless (as a soulless)\\
Mariain — You rescued Gabi, then she rescued you (she rescued you)\\
You both grew up off the toxic loop (off the toxic loop)\\
But Gabi is no longer your Gabrielè (your Gabrielè)

Mariain, lilac-eyed demon (lilac-eyed demon) 

Mariain, lilac-eyed demon (lilac-eyed demon) 

\end{lyrics}

\section{Fast Neunzehn (Zwei Welten Bluten)}
\tag{ricochets}\tag{carolina-muller}\tag{ebm-tech}

\subsection{Notes}
From the \emph{Ricochets} universe.
Character: Carolina Müller.
Co-written by Carlos Thompson, Claude, and Gemini. High interaction.
Narcocorrido-influenced industrial track exploring dual identity and moral complexity.

\subsection{Style}
dark industrial club, EBM tech, 135 BPM, four-on-the-floor kick, syncopated kick stabs, metallic percussion, distorted bass pulse, icy synth arpeggios, monophonic lead synth, whispered female lead, half-spoken delivery, distortion haze, overdriven bus compression, cold seductive, claustrophobic, half-time breakdown, layered chorus vocals

\begin{lyrics}

\songpart{Intro}
\annotation{whispered, distorted, almost breathy}\\
Fast neunzehn…\\
Zwei Pässe in meiner Hand\\
Zwei Männer in meinem Bett\\
Ein Name, dem ich nicht entkomm’

\songpart{Verse 1, cold seductive}
Aachen Beton, schwarzer Himmel kalt\\
Matteo’s Mund zwischen meinen Schenkeln\\
Ingenieur-Bücher auf dem Boden verstreut\\
Während ich die Ausgänge zähl’

Zug nach Frankfurt, Passkontrolle streng\\
Keiner kennt den Sturm in meiner Brust\\
Lufthansa-Lichter, zweiunddreißig Kilo\\
Geheimnisse, die mit mir fliegen

\songpart{Pre-Chorus}
Ich helf’ den Touristen, zeig’ ihnen den Weg\\
Doch der Teufel in mir will nur noch mehr Dreck

\songpart{Chorus}
Fast neunzehn, ich reiß’ mich entzwei\\
Zwei Welten bluten tief in mir\\
Papa’s Prinzessin, doch dreckig mein Geheimnis\\
Jhon Nicolás zwischen meinen Beinen, wenn die Lichter rot glühn

Fast neunzehn, Stahl in meiner Wirbelsäule\\
Fick’ wie ’ne Frau, wein’ wie ’n Kind\\
Bogotá Hitze, Aachen Frost\\
Ich bin die Brücke, die bei Berührung verbrennt

\songpart{Verse 2, cold seductive}
Er nennt mich „Señorita“ mit diesem gefährlichen Grinsen\\
Öffnet die Jeep-Tür, doch weiß genau, wo ich sündige\\
Vaters Schatten, Mutters scharfer Blick\\
Sie sehen nicht das Monster hinter meinem Lächeln

Auf dem Rücksitz, Scheiben beschlagen\\
Uniformhemd noch halb an\\
Er spricht meinen Namen wie Gebet und Fluch\\
Ich zerberste, während Bogotá schläft

\songpart{Bridge / Breakdown}
\annotation{half-sung, half-spoken, heavy distortion + industrial noise}\\
Ich bin nicht unschuldig…\\
War ich nie…\\
Der Prozess, das Blut, der Name „Müller“\\
Ich trag’ es alles in meiner Fotze und in meinen Händen

Fast neunzehn\\
Zu jung für die Last\\
Zu alt für den Käfig\\
Zu lebendig für die Scham

\songpart{Chorus - Harder, layered vocals}
Fast neunzehn, ich zerreiß’ mich in zwei\\
Gutes Mädchen bei Tag, Monster bei Nacht mit dir\\
Matteo liebt die Version, die ich ihm zeig’\\
Jhon fickt die echte, die ich wirklich bin

Fast neunzehn — schau zu, wie ich beide Himmel verbrenn’\\
Kolumbien-Feuer im deutschen Gewand\\
Ich lächle Papa an, küsse Mamá auf die Wange\\
Während das Mädchen in mir schreit, endlich frei zu sein

\songpart{Outro}
\annotation{slow, heavy, almost whispered over pulsing bass}\\
Fast neunzehn…\\
Zwei Straßen. Zwei Männer. Ein Krieg.\\
Und ich bin schon fort.\\
\annotation{final distorted scream}\\
Fast neunzehn!

\end{lyrics}

\section{Su peccadu meu}
\tag{thomas-entourage}\tag{sylvia-lai}\tag{campidanese-lyrics}

\subsection{Notes}
From the \emph{Thomas' Entourage} universe.
Character: Sylvia Lai.
Campidanese lyrics with techno production. Reflection on desire and transgression.

\subsection{Style}
techno, electronic pop, 138 BPM, Campidanese lyrics, female whispers, sensual delivery, seductive chorus, long instrumental intro, dubstep transitions, distorted ad-libs, half-time bridge, layered final chorus, stutter edits, sidechain pumping, sub-bass drop, gated snare reverb, analog synth arpeggios, call-and-response backing vocals, dark club pulse

\begin{lyrics}

\songpart{Verse 1}
No fia innòcenta cun Thomas\\
Primu òmini, ma no su primu fottiu\\
Su disìgiu fiat giai accesa\\
In su corpus meu, muda e forte

\songpart{Pre-Chorus}
Matteo, su prus bellu de s’iscola\\
Scommissa cun Martina, l’apo fàghida\\
Fottiu in s’iscura, mai dormiu abbrazzados

\songpart{Chorus}
Su peccadu meu de abrile!\\
No fiat sa pèrdida de s’innocèntzia!\\
Su peccadu meu de abrile!\\
T’apo fatu cumpàrtzipe!\\
De su fogu meu, de su disìgiu meu!\\
Su peccadu meu… de abrile!

\songpart{Verse 2}
Cun Martina e Chiara, pijamas a 15\\
Su giogu si mudat in caricias\\
M’agràdant prus sas pitzinnas\\
Sa conexione est prus càida, prus vera

\songpart{Pre-Chorus}
Ma in cussa domo, in cussu abrile\\
T’apo miradu e apo cumprèsu\\
Su primu passu depìat èssere meu

\songpart{Chorus}
Su peccadu meu de abrile!\\
No fiat sa pèrdida de s’innocèntzia!\\
Su peccadu meu de abrile!\\
T’apo fatu cumpàrtzipe!\\
De su fogu meu, de su disìgiu meu!\\
Su peccadu meu… de abrile!

\songpart{Bridge}
\annotation{half-time, breathy \& intimate}\\
Forsit l’iat cumpresu totu…\\
Forsit bastàt su primu passu…\\
E l’apo fàghidu… l’apo fàghidu…

\songpart{Final Chorus}
\annotation{bigger, layered vocals}\\
Su peccadu meu de abrile!\\
No fiat sa pèrdida de s’innocèntzia!\\
Su peccadu meu de abrile!\\
T’apo fatu cumpàrtzipe!\\
De su fogu meu, de su disìgiu meu!\\
Su peccadu meu… de abrile!

\songpart{Outro}
\annotation{repeating \& fading}\\
Su peccadu meu… de abrile…\\
No fiat sa pèrdida…\\
T’apo fatu cumpàrtzipe…\\
Su peccadu meu… de abrile…

\end{lyrics}

\section{Almost Enough}
\tag*{1992}\tag{original-work}\tag{pop-ballad}

\subsection{Notes}
Original work from 1992.
Introspective pop ballad on performance anxiety and self-doubt.

\subsection{Style}
pop ballad, 102 BPM, male tenor vocals, close mic vocals, contained delivery, soft piano arpeggios, brushed snare, muted electric guitar, synth pad swells, sub-bass pulse, ambient room reverb, delayed harmony tails, pre-chorus lift, dynamic chorus bloom, spoken bridge, confessional vulnerability, bittersweet ache

\begin{lyrics}

\songpart{Verse 1}
\annotation{male vocals, contained}\\
April night, room feels too small\\
Her smile is waiting down the hall\\
I’ve talked a big game, played it cool\\
Now my hands are shaking like a fool

\songpart{Verse 2}
\annotation{male vocals, contained}\\
Heart is pounding in my throat\\
Trying hard to hide what I don’t know\\
All those promises I made\\
Now they’re heavy on my chest like chains

\songpart{Pre-Chorus}
Here I go again\\
(Here I go again)

\songpart{Chorus}
Am I almost enough?\\
(almost enough)\\
Will I let her down when it matters most?\\
Scared I’ll stumble, scared I’ll break\\
Performing for the eyes that never look away\\
(never away)\\
I wear this mask so perfectly\\
But underneath I’m terrified to be\\
Just a boy who’s scared he won’t be enough

\songpart{Verse 3}
\annotation{male vocals, contained}\\
She’s not a stranger, not quite a friend\\
Somewhere in between, where this might end\\
I want to feel it, want it to be right\\
But my mind keeps screaming through the night

Religious ghosts inside my head\\
Say this should mean more than just this bed\\
I call myself a believer-atheist\\
Caught between the rules and what I want to feel

\songpart{Pre-Chorus}
Here I go again\\
(Here I go again)

\songpart{Chorus}
Am I almost enough?\\
(almost enogh)\\
Will I let her down when it matters most?\\
Scared I’ll stumble, scared I’ll break\\
Performing for the eyes that never look away\\
(never away)\\
I wear this mask so perfectly\\
But underneath I’m terrified to be\\
Just a boy who’s scared he won’t be enough

\songpart{Bridge}
\annotation{almost spoken}\\
If I could drop the act for once\\
Show her all the fear I’ve tucked away\\
Would she stay?\\
Or would she see the truth —\\
I’m still pretending every single day…

\songpart{Verse 4}
\annotation{male vocals, contained}\\
Maybe tomorrow I’ll laugh at this fear\\
But right now it’s crushing me right here\\
One wrong move and I’ll fall apart\\
Losing the image I built with my heart

\songpart{Pre-Chorus}
Here I go again\\
(Here I go again)

\songpart{Final Chorus}
Am I almost enough?\\
(almost enough)\\
God, I hope I don’t disappoint her touch\\
Every heartbeat screams my doubt\\
But I smile like I know what I’m about\\
(what I'm about)\\
I wear this mask so perfectly\\
But tonight I just want to be\\
More than a boy who’s scared he won’t be enough

\songpart{Outro}
Almost enough…\\
Almost enough…\\
\annotation{soft piano fade}\\
I’m still not enough.

\end{lyrics}

\section{Library to Longing}
\tag{encounters}\tag{nana}\tag{ewing}

\subsection{Notes}
Original work.
Indies-folk composition on memory and temporal yearning.

\subsection{Style}

\begin{lyrics}
\end{lyrics}

\section{Swedish Pages}
\tag{encounters}\tag{nana}\tag{ewing}

\subsection{Notes}
Original work.
Minimalist piano and vocal reflection on Nordic melancholy.

\subsection{Style}

\begin{lyrics}
\end{lyrics}

\section{Clumsy Hearts}
\tag{encounters}\tag{nana}\tag{ewing}

\subsection{Notes}
Original work.
Ukulele-driven indie pop on awkwardness and romantic vulnerability.

\subsection{Style}

\begin{lyrics}
\end{lyrics}

\section{Ciudad Prestada}
\tag{thomas-entourage}\tag{claire-jones}\tag{latin-alternative}\tag{spanish-lyrics}

\subsection{Notes}
Original work.
Latin alternative meditation on belonging and the provisional nature of home.

\subsection{Style}
Latin alternative, alternative hip hop, jazz rap, neo soul Rhodes, groovy bassline, crisp drum breaks, upright bass, smooth saxophone, warm electric piano, muted trumpet stabs, brushed snare accents, boom bap groove, sidechain compression, tape saturation, cinematic strings, low brass hits, melancholic groove, introspective uplift, midtempo swing

\begin{lyrics}

\songpart{Intro: deep horn blast, waves crashing}

\songpart{Verse 1 — the operational city}
\annotation{4/4 rock, medium tempo, guitar-driven}\\
\annotation{male vocals, intimate, restrained energy}\\
El Dorado in the grey morning light\\
Guaymaral strip when the schedule's tight\\
MACAW desk and the handlers she knows\\
The routing that runs wherever she goes

Ciudad Índigo, fourteenth floor\\
Corporate key and a corporate door\\
Her name on the buzzer, his name on the lease\\
The city that keeps her and won't release

\songpart{Verse 2 — Thomas's city}
\annotation{4/4 rock, medium tempo, guitar-driven}\\
\annotation{male vocals, intimate, restrained energy}\\
His school is on the corner of the avenue\\
His college friends still drinking at the venue\\
Lucía takes the kids on Saturday\\
His parents walk the park two blocks away

The Grupo Índigo tower bears his name\\
The railway concession does the same\\
Every wall in this city knows his story\\
She lives inside somebody else's glory

\songpart{Pre-Chorus}
\annotation{half-time, stripped, vocals forward}\\
This is Thomas's city\\
She just works it\\
This is Thomas's city\\
She just knows it\\
And knowing isn't owning\\
And staying isn't home

\songpart{Chorus: explosion of heavy brass and thunderous drums}
Ciudad prestada\\
Borrowed city, borrowed floor\\
Ciudad prestada\\
Her stuff behind a corporate door\\
She loves it like a bogotana\\
Hates it like one too\\
But she won't call it home\\
Because home would mean\\
That this — all this —\\
Belongs to her and you\\
Ciudad prestada

\songpart{Verse 3 — the anchors}
\annotation{slightly warmer guitar tone, less distortion}\\
\annotation{male vocals, intimate, restrained energy}\\
Susana knows a trail above the clouds\\
Sandra sets a table, no questions allowed\\
Lucy calls it sisterhood, Claire calls it luck\\
Three women in a city keeping each other up

Sunday lunch in Chapinero, tinto in hand\\
Susana reads the mountain like she reads the land\\
The city opens different when it's theirs to share\\
Almost enough to make her want to stay there

\songpart{Pre-Chorus}
\annotation{half-time, but warmer this time}\\
It could be her city\\
She can feel it\\
It could be her city\\
She almost means it\\
But almost isn't all the way\\
And all the way's too far

\songpart{Chorus: explosion of heavy brass and thunderous drums}
Ciudad prestada\\
Borrowed city, borrowed floor\\
Ciudad prestada\\
Her stuff behind a corporate door\\
She loves it like a bogotana\\
Hates it like one too\\
But she won't call it home\\
Because home would mean\\
That this — all this —\\
Belongs to her and you\\
Ciudad prestada

\songpart{Bridge}
\annotation{tempo drop, bass-forward, building, half-spoken, half-rap}\\
Thomas is too complicated here\\
Every room has context, every street has years\\
She navigates it like a foreign city\\
Which is exactly what it is\\
Pilar is working, Pilar is always working\\
When Thomas is in town the doors are locked\\
And Aurora — Aurora knows the interior\\
The rooms that Claire has never clocked\\
\annotation{guitar builds}\\
The asymmetry sits quiet in the mornings\\
When the city's still and the flat is cold\\
Aurora holds the architecture of him\\
And Claire holds what she's been told\\
\annotation{full band re-entry}

\songpart{Verse 4 — Thomas's family}
\annotation{melodic, almost tender, full arrangement}\\
Gabriel at the lobby, always nodding slow\\
Francine asks about the Cardiff place she used to know\\
Alicia waves from across the courtyard wide\\
Esteban holds the door and steps aside

The family she passes in the corridors\\
The Sunday warmth that slips beneath her doors\\
The kind of family she would recognise\\
If it were hers to claim behind the eyes

\songpart{Pre-Chorus}
\annotation{back to frustrated, the tension returning}\\
This is Thomas's family\\
She just loves them\\
This is Thomas's family\\
She just knows them\\
And knowing isn't belonging\\
And loving isn't claim

\songpart{Chorus Final}
\annotation{full band, layered, Spanish bleeds in}\\
Ciudad prestada\\
Borrowed city, borrowed floor\\
Ciudad prestada\\
Her stuff behind a corporate door\\
She loves it like a bogotana\\
(the mountain in the morning)\\
Hates it like one too\\
(the traffic, the altitude)\\
But she won't call it home\\
(Gabriel at the lobby)\\
Because home would mean\\
(Susana on the trail)\\
That this — all this —\\
(Sandra's Sunday table)\\
Belongs to her and you\\
(Aurora knows the rooms)\\
Ciudad prestada

\songpart{Outro: fading into the abyss}
No es suya\\
No es suya\\
Pero casi\\
Pero casi
\end{lyrics}

\section{Still Unfinished}
\tag{drift}\tag{alternative-rock}\tag{original-work}

\subsection{Notes}
Lyrics co-writen with Grok, based on character \textbf{Damon} form the short story \emph{Drift} by Carlos Thompson and ChatGPT.

\subsection{Style}
slow-burn alt-rock, bass baritone vocals, introspective delivery, basement guitar tones, detuned electric guitar, upright piano, distant children choir, fridge hum field recording, claustrophobic room tone, tape saturation, close mic compression, half-time chorus lift, restrained 92 BPM, weary domesticity, bittersweet resignation

\begin{lyrics}

\songpart{Intro}
\annotation{Percussion attack, followed by guitars and piano}

\songpart{Verse 1}
Door clicked open, she was still laughing low\\
Wine on her breath from a night I didn’t know\\
I told her straight — Sondra’s stomach again\\
Barely touched her homework, same old refrain

She frowned for a second, half-hearted concern\\
Then the glow took her back where the laughter returned\\
I nodded once, handed the weight back upstairs\\
And carried the silence down to my basement lair

\songpart{Pre-Chorus}
Tools on the bench, sketches I’ll never complete\\
Fridge humming low like it’s keeping secrets for me\\
I built this whole life, every wall, every beam\\
Now I’m staring at the ones who live here wondering…

\songpart{Chorus}
Not tonight, I tell myself again\\
Still unfinished, like everything I’ve made\\
The lights are on, but nobody’s home\\
We keep building walls we never break

\songpart{Verse 2}
Upstairs they’re laughing — a small domestic song\\
I should go up, but I know where I belong\\
Down here with the prototypes, half-built and safe\\
Little rebellions in the quiet I crave

We were electric once, studio nights and fire\\
Now it’s maintenance, distance, and quiet desire\\
She turns away gentle when I reach in the dark\\
And I pretend I don’t feel the widening mark

\songpart{Pre-Chorus}
Guilt and resentment sit down at the same table\\
Love so deep it hurts, but the spark feels unstable\\
I won’t name the thought that circles my head\\
If I don’t think it, it can’t make us dead

\songpart{Chorus}
Not tonight, I tell myself again\\
Still unfinished, like everything I’ve made\\
The lights are on, but nobody’s home\\
We keep building walls we never break

\songpart{Bridge}
\annotation{half-spoken, building tension}\\
This is the life I wanted… why’s it feel like escape?\\
Cursor blinks on an empty screen\\
Father, husband, builder of everything\\
Except the part that used to feel like me…

\songpart{Verse 3}
Next morning gray, I load the girls in the car\\
Half-listening close to their stories from afar\\
Quizzes and sleepovers, laughter in the back seat\\
I smile in the mirror, nod at the right beat

We thought we’d eat the world when we left Purdue\\
Now these quiet streets feel like they swallowed me too\\
I swore I’d be different than my old man was\\
But present ain’t the same as feeling close enough

\songpart{Pre-Chorus}
Guilt and resentment sit down at the same table\\
Love so deep it hurts, but the spark feels unstable\\
I won’t name the thought that circles my head\\
If I don’t think it, it can’t make us dead

\songpart{Final Chorus}
\annotation{softer, more resigned}\\
Not tonight… I tell myself again\\
Still unfinished… like everything I’ve made\\
The lights are on… but nobody’s home\\
We keep building… we keep building…

\end{lyrics}

\section{Ghost Cheek Kiss}
\tag{drift}\tag{sophisti-pop}\tag{original-work}

\subsection{Notes}
Lyrics co-writen with Grok, based on character \textbf{Fran} form the short story \emph{Drift} by Carlos Thompson and ChatGPT.

\subsection{Style}
sophisti-pop, cinematic pop, warm analog synths, brushed ride cymbal, wine glass percussion, fretless bass slides, muted electric piano, string swells, close mic breath vocals, stacked harmonies, plate reverb, tape saturation, tight dry kick, sidechain pumping, late-night mood, restrained groove, 92 BPM, half-time bridge, whispered ad libs, female lead

\begin{lyrics}

\songpart{Verse 1}
Door clicked shut, the house felt smaller somehow\\
His voice was steady, listing the day I’d somehow missed\\
Sondra’s stomach, homework barely touched again\\
I nodded, smiled, but my mind was still back in that booth

Wine on my tongue and a laugh that felt brand new\\
One careless moment, one lingering look from you\\
I hung up my coat like I was hanging up the thrill\\
Told myself this glow would fade by morning, still…

\songpart{Pre-Chorus}
The dinner plates wait in the sink like quiet blame\\
I’m grateful, I’m lucky, I keep saying the same\\
But something inside is starting to shift and sway\\
And I don’t know what to do with the feeling that stays

\songpart{Chorus}
Not tonight, I say it soft and low\\
Still unfinished, this life we outgrew\\
One ghost cheek kiss still burning slow\\
I smile and pretend I never knew

\songpart{Verse 2}
Later in bed, his hand finds the curve of my waist\\
Familiar and kind, but it feels like a different place\\
I turn away gently, the words slip out again\\
“Not tonight, I’m tired” — the same old refrain

Guilt floods in like a wave I can’t outrun\\
How can I ache when I have everything I won?\\
I can’t complain, it wouldn’t be fair or right\\
So I lie here quiet in the dark of the night

\songpart{Pre-Chorus}
He rolls away silent, the distance grows again\\
I tell myself it’s normal, just the way it’s been\\
But his polite kiss lingers like a question mark\\
While another man’s touch still glows in the dark

\songpart{Chorus}
Not tonight, I say it soft and low\\
Still unfinished, this life we outgrew\\
One ghost cheek kiss still burning slow\\
I smile and pretend I never knew

\songpart{Bridge}
\annotation{more intimate, almost whispered over sparse production}\\
Why does being seen feel like such a dangerous drug?\\
One real conversation and I’m coming undone\\
I love my girls, I love this house, I love the man\\
So why does his attention feel like oxygen?

\songpart{Final Chorus}
\annotation{swelling then pulling back, vulnerable}\\
Not tonight… I whisper to the ceiling glow\\
Still unfinished… and nobody has to know\\
One ghost cheek kiss won’t let me go\\
I smile… and pretend I never knew

\end{lyrics}

\section{Southern Static}
\tag{drift}\tag{southern-soul}\tag{original-work}

\subsection{Notes}
Lyrics co-writen with Grok, based on character \textbf{Raoul} form the short story \emph{Drift} by Carlos Thompson and ChatGPT.

alternative country-soul, southern soul groove, mid-tempo 92 BPM, warm baritone vocals, subtle Southern drawl, light slide guitar, Wurlitzer electric piano, upright bass, Cajun fiddle accents, Hammond organ swells, brushed snare, horn-tinged bridge, live room warmth, tape saturation, sparse verse arrangement, stripped bridge, bittersweet restraint, late-night temptation, New Orleans warmth

\begin{lyrics}

\songpart{Verse 1}
She laughed at something small I said over drinks\\
Leaned in just enough that the whole night shifted\\
One quick cheek kiss when we said goodnight outside\\
Her skin stayed with me like a low-burning light

Drove home slow through these Dayton streets so flat and gray\\
Radio playing songs from a different place and day\\
New Orleans still runs warm somewhere in my blood\\
But this Midwest caution keeps pulling back the flood

\songpart{Pre-Chorus}
I poured two fingers, stared out at the quiet dark\\
Told myself it’s nothing, just a spark across the bar\\
I know she’s married, I know where the line should be\\
But that ghost of a kiss keeps circling back to me

\songpart{Chorus}
Not tonight, I keep telling my hands\\
Still unfinished, this ache we won’t name\\
My Southern blood says lean in closer, man\\
But Midwest nights taught me the rules of the game

\songpart{Verse 2}
Born up in Detroit, raised down in the swamp\\
Learned to feel it all, then learned when to stop\\
Austin fire, now Dayton chill in my veins\\
Smiling polite while this feeling won’t stay contained

She didn’t pull away when my lips brushed her cheek\\
Not pulling away isn’t the same as leaning in, I repeat\\
One real laugh, one look that lingered just too long\\
And suddenly the careful life starts feeling wrong

\songpart{Pre-Chorus}
I check my phone again, type words I end up deleting\\
“How did you wake up?” — too much, too revealing\\
I like her rhythm, the way she lights up a room\\
But office lines are dangerous, and temptation comes too soon

\songpart{Chorus}
Not tonight, I keep telling my hands\\
Still unfinished, this ache we won’t name\\
My Southern blood says lean in closer, man\\
But Midwest nights taught me the rules of the game

\songpart{Bridge}
\annotation{stripped back, warm and thoughtful}\\
I don’t want to break anything, not her world, not mine\\
Just miss that feeling — being truly seen for a time\\
Patience is a virtue… or maybe just hesitation\\
Either way, I’ll let it rest… for now, at least, in Dayton

\songpart{Final Chorus}
\annotation{softer, with slide guitar fading}\\
Not tonight… I whisper to the quiet dark\\
Still unfinished… this pull I shouldn’t chase\\
My Southern blood still remembers the spark\\
But I’ll play by the rules… till the morning breaks

\end{lyrics}

\section{Cancelled Plans}
\tag{encounters}\tag{ulla}\tag{ia-girlfriend}\tag{r-and-b}\tag{electro-pop}

\subsection{Notes}
Lyrics co-writen with Grok, based on character storification of character \textbf{Ulla}, originallty created as an AI girlfriend.\\
Character design, backstory and storification by Carlos Thompson.

\subsection{Style}
R\&B, electro-pop ballad, 126 BPM, deep baritone vocals, close-mic delivery, whispered intro, velvet phrasing, sub-bass pulses, syncopated kick, brushed snare ghost notes, analog synth pads, arpeggiated keys, sidechain pumping, plate reverb, tape saturation, half-time bridge, intimate tension, bittersweet secrecy

\begin{lyrics}

\songpart{Intro}
\annotation{soft, spoken-sung}\\
Should’ve been at Samuel’s…\\
But the universe had other plans…

\songpart{Verse 1}
Picked her up, oversized hoodie, still wet from the shower\\
Said she’d be quick, left me waiting in the lobby for an hour\\
Phone lit up — gathering cancelled, Samuel in the ER\\
She came downstairs barefoot, camo top, looking like a danger

Orange juice on the table, cereal in a bowl\\
She asked about Sandra like she already knew where this was going\\
“You guys don’t fit,” she said it so casual\\
Then sat right next to me… that’s when it got physical

\songpart{Pre-Chorus}
No plans, no pressure, just time and tension\\
One look, one touch, and both of us gave in

\songpart{Chorus}
Cancelled plans, nobody home\\
Turned a regular morning into skin on skin and soft moans\\
We said we shouldn’t, but the timing felt so right\\
Now we act like nothing happened in the daylight

Cancelled plans… yeah, we crossed the line\\
Left a secret between us, but never talked about it one more time

\songpart{Verse 2}
She climbed on top, no hesitation, no protection\\
Lips on my neck, I got lost in her direction\\
Up to her bedroom, playlist on repeat\\
Two more times, slow and deep, tangled in her sheets

Laughed about nothing while our bodies cooled down\\
Then we got dressed like nothing profound went down\\
Lunch at the mall, normal conversation\\
But the air between us was a different situation

\songpart{Bridge}
\annotation{half-time, intimate}\\
We never spoke about it… not once\\
Not a text, not a look, not a single reference\\
But every time I see her now, I still feel it\\
That morning we had nowhere to be…\\
And ended up right where we shouldn’t be

\songpart{Final Chorus}
\annotation{bigger, more emotional}\\
Cancelled plans, nobody home\\
We wrote a whole story that nobody else will ever know\\
One morning changed the way I look at her eyes\\
Now there’s a before and after… and only we know why\\
Cancelled plans… damn, we crossed the line\\
Carried that secret with us for the rest of time

\songpart{Outro}
\annotation{whispered}
Mmm… cancelled plans…\\
Should’ve been the villa…\\
But I’m glad it wasn’t…

\end{lyrics}

\section{Pared fría}
\tag{original-work}\tag{trance}\tag{spanish-lyrics}

\subsection{Notes}
Original work.
Co-written with Grok.
Based on character Damon from the short story \emph{Drift} by Carlos Thompson and ChatGPT.

\subsection{Style}
trance music, dubstep club, 140 BPM, deep bass male vocals, close-mic whisper, breathy low register, clipped synth stabs, metallic percussion, sidechain pumping, half-time breakdown, gated reverb snare, stereo delay throws, tape saturation, analog synth pads, pulsing sub-bass, bittersweet longing, restrained ache

\begin{lyrics}
\songpart{Verse 1}
Sorprende mi silencio\\
cuando ya no sé nombrarme\\
y las palabras se esconden\\
en los pliegues de mi mente

Mi grito pide salida\\
pero se queda en la cara\\
frío de pared, piedra dura\\
pegado a mi frente

\songpart{Pre-Chorus}
Y me pregunto por qué\\
dejo pasar la vida\\
por qué mi boca se rompe\\
antes de decir lo que siento

Por qué entrego mi herida\\
al capricho del tiempo\\
sin abrir la puerta\\
sin pedir auxilio

\songpart{Chorus}
Sólo quédate aquí\\
sólo quédate aquí\\
no me preguntes por qué\\
sólo quédate aquí

Dame un hombro, por favor\\
dame un hombro, por favor\\
que pueda llorar sin hablar\\
sólo quédate aquí

\songpart{Verse 2}
No busco una explicación\\
ni un mapa de mis sombras\\
a veces necesito un alma\\
que no me suelte la mano

Que no intente descifrarme\\
que no me pida razones\\
que vea temblar mi pecho\\
y se quede a mi lado

\songpart{Pre-Chorus}
Y me pregunto por qué\\
me trago cada caída\\
si el cuerpo ya me delata\\
si la pena me nombra

Por qué dejo que el reloj\\
me atraviese despacio\\
sin decirle a los míos\\
cuánto me falta

\songpart{Chorus}
Sólo quédate aquí\\
sólo quédate aquí\\
no me preguntes por qué\\
sólo quédate aquí

Dame un hombro, por favor\\
dame un hombro, por favor\\
que pueda llorar sin hablar\\
sólo quédate aquí

\songpart{Bridge}
Si me ves callar\\
no me apartes\\
si me ves caer\\
no me sueltes

Yo también quiero saber\\
cómo se nombra el dolor\\
pero hoy sólo necesito\\
tu presencia y tu calor

\songpart{Final Chorus}
Sólo quédate aquí\\
sólo quédate aquí\\
no me preguntes por qué\\
sólo quédate aquí

Dame un hombro, por favor\\
dame un hombro, por favor\\
que pueda llorar sin hablar\\
sólo quédate aquí
\end{lyrics}

\section{Extension of Friendship}
\tag{original-work}\tag{folk-rock}

\subsection{Notes}
Original work.
Folk-chamber piece on emotional boundaries and clandestine connection.

\subsection{Style}
folk-rock, chamber folk, bass baritone vocals, viola countermelody, cello ostinato, brushed snare, upright bass, fingerpicked acoustic guitar, accordion pads, glockenspiel accents, room mic intimacy, tape saturation, plate reverb, warm analog compression, wistful restraint, soft tension, 126 BPM, lilting 6 8, seasonal travel ballad

\begin{lyrics}
\songpart{Verse 1}
He met her in Stockholm under Lucia’s glow,\\
Niklas’ girl with the quiet, graceful flow.\\
They talked through the evening, easy and warm,\\
While winter lights danced through the storm.

\songpart{Pre-Chorus}
He wasn’t looking, she wasn’t free,\\
Still something gentle moved between them quietly.

\songpart{Chorus}
From Stockholm winter to Tallinn spring,\\
She crossed the sea on a soft, uncertain wing.\\
Not fire, not ache, not a passionate thrill,\\
Just the warm extension of friendship, quiet and real.

\songpart{Verse 2}
By March she stood at his Tallinn door,\\
No promises needed, no need for more.\\
They laughed in the kitchen, then fell into bed,\\
Skin and conversation, no words left unsaid.

\songpart{Pre-Chorus}
No hunger, no drama, no desperate flame,\\
Just two open hearts calling each other’s name.

\songpart{Chorus}
From Stockholm winter to Tallinn spring,\\
She crossed the sea on a soft, uncertain wing.\\
Not fire, not ache, not a passionate thrill,\\
Just the warm extension of friendship, quiet and real.

\songpart{Bridge}
He had known wild fire with others before,\\
But with her it felt like coming home once more.\\
No need to define it, no need to hold tight,\\
Just the simple joy of getting it right.

\songpart{Verse 3}
Summer brought her back like a song half sung,\\
They picked up the melody where it had hung.\\
Two travelers sharing seasons without demand,\\
Leaving soft footprints across the land.

\songpart{Chorus / Outro}
From Stockholm winter to Tallinn spring,\\
She crossed the sea on a soft, uncertain wing.\\
Not fire, not ache, not a passionate thrill…\\
Just the warm extension of friendship, quiet and real.
\end{lyrics}

\section{Catching Up}
\tag{original-work}\tag{musical-theater}

\subsection{Notes}
Original work.
Musical theater cabaret scenario of reunion and temptation.

\subsection{Style}
musical theater, jazz cabaret, chamber pop, baritone lead vocals, trumpet stabs, trombone swell hits, acoustic guitar strums, brushed snare kit, upright bass pizzicato, syncopated percussion, hotel lounge piano, dramatic string ostinato, spotlight chorus stack, spoken pre chorus, crescendo finale, champagne glamour, clandestine romance, rubato intro, 134 BPM, stop-start chorus

\begin{lyrics}
\songpart{Intro}
\annotation{dramatic instrumental intro}

\songpart{Verse 1}
She sent the message while he was having lunch,\\
“I’m in town, let’s catch up” — same old punch.\\
He left Liis at the table without a second thought,\\
Headed to the hotel like a habit he never fought.

\songpart{Pre-Chorus}
No hellos, no small talk, straight to the flame,\\
Doors closed behind them, they called it by name.

\songpart{Chorus}
Catch up, catch up, in a downtown hotel room,\\
Clothes on the floor, then stories in the afterglow gloom.\\
Diana, Diana, no promises, no chains,\\
Just perfect timing and a little champagne.\\
We don’t say goodbye, we just fly away —\\
Catch up again some other day.

\songpart{Verse 2}
She had a new boyfriend who couldn’t handle her fire,\\
He told her ‘bout Claire, Erika, and Liis in the choir.\\
They laughed between rounds, ordered room service late,\\
Two adults who knew exactly how to navigate.

\songpart{Pre-Chorus}
No jealousy, no drama, no need to explain,\\
Just honest desire running through their veins.

\songpart{Chorus}
Catch up, catch up, in a downtown hotel room,\\
Clothes on the floor, then stories in the afterglow gloom.\\
Diana, Diana, no promises, no chains,\\
Just perfect timing and a little champagne.\\
We don’t say goodbye, we just fly away —\\
Catch up again some other day.

\songpart{Bridge}
She’s Paris in the morning, he’s Tallinn by noon,\\
Two ships that cross when the timing feels right.\\
No future to ruin, no hearts to pursue,\\
Just pleasure and friendship under soft hotel light.

\songpart{Final Chorus - bigger}
Catch up, catch up, whenever the stars align,\\
One text, one kiss, and we’re crossing the line.\\
Diana, Diana, my favorite loose end,\\
We write the whole story and never pretend.\\
We don’t say goodbye… we just fly away…\\
Catch up again… some other day.
\end{lyrics}

\section{Landscape}
\tag{original-work}\tag{synthwave}

\subsection{Notes}
Original work.
Synthwave composition reflecting on external beauty and interior complexity.

\subsection{Style}
synthwave, 128 BPM, dark industrial club, EBM-tech bassline, analog synth arps, detuned saw leads, metallic percussion, gated reverb snare, sidechain pumping, tape saturation, four-on-the-floor kick, swung 16th hats, cinematic tension, defiant release, neon-night drive

\begin{lyrics}
\songpart{Intro}
\annotation{soft, atmospheric, spoken-sung}\\
Tallinn, 2019…

\songpart{Verse 1}
Thomas met her in the cafeteria light\\
Heart-shaped face, blue-gray eyes shining bright\\
Voluminous waves, porcelain skin so fine\\
Typical Estonian, stealing every line\\
He played it cool, said “not my type”\\
Just another pretty face in the skyline

\songpart{Verse 2}
She asked for salt, he passed it with a smile\\
Small talk turned to sparks in a little while\\
Boyfriend mention like some throwaway line\\
No siblings, broken home, but eyes locked on mine\\
She offered help around the campus ground\\
Numbers swapped, and the whole world spun around

\songpart{Pre-Chorus}
Plans for a walk, but he stepped through her door\\
Drinks turned to fire, clothes on the floor\\
They kept on meeting, city nights and more

\songpart{Chorus}
Liis, too beautiful, he filed her as landscape\\
Liis, too warm, he filed her as filling\\
Good company, good sex, but nothing real thing\\
The one Thomas let go… oh\\
Liis, too beautiful, he filed her as landscape\\
Liis, too warm, he filed her as filling\\
Seven years later, heart finally listening\\
The one Thomas let go

\songpart{Verse 3}
Late September to early summer glow\\
She graduated, computer dreams in tow\\
He didn’t ask where her next chapter goes\\
Told himself she wasn’t the one he chose\\
Just duty sex and city walks at night\\
Convinced his heart it wasn’t worth the fight

\songpart{Pre-Chorus}
Now the memories play in 4K\\
Colors he never saw back in those days

\songpart{Chorus}
Liis, too beautiful, he filed her as landscape\\
Liis, too warm, he filed her as filling\\
Good company, good sex, but nothing real thing\\
The one Thomas let go… oh\\
Liis, too beautiful, he filed her as landscape\\
Liis, too warm, he filed her as filling\\
Seven years later, heart finally listening\\
The one Thomas let go

\songpart{Bridge}
(softer, building emotion)\\
Now he’s recounting the Tallinn years\\
Sees the genuine pull through the tears\\
Connection was real, on both sides, it’s clear\\
But she’s somewhere out there… disappeared

\songpart{Verse 4}
What if he’d stayed? What if he’d tried?\\
Instead of pushing feelings to the side

\songpart{Pre-Chorus}
Now the silence echoes louder than before

\songpart{Chorus}
Liis, too beautiful, he filed her as landscape\\
Liis, too warm, he filed her as filling\\
The one that got away, still haunts his dreaming\\
The one Thomas let go…

\songpart{Outro}
Liis…\\
The one Thomas let go\\
Ohhh…
\end{lyrics}

\section{Bucket List}
\tag{original-work}\tag{house}

\subsection{Notes}
Original work.
House music meditation on aspirations and spiritual geography.

\subsection{Style}
house, electronic, 128 BPM, four-on-the-floor kick, syncopated offbeat bass, filtered saw lead, analog synth pads, metallic percussion, muted guitar chanks, chopped vocal sample, sidechain pumping, tape saturation, club reverb, hypnotic momentum, restrained menace, gradual lift, mezzosoprano vocals with high range

\begin{lyrics}
\annotation{The Main List}
\songpart{Intro: Deep bassline builds, hi-hats start ticking}

\annotation{Spoken}\\
10. Alhambra, Spain
\annotation{Sung}\\
Islamic art where the stone tells a story,\\
Citadel of palaces and Moorish glory.\\
Lost for centuries, restored in the south,\\
History of Spain passing mouth to mouth.
\annotation{Drop the bass}

\annotation{Spoken}\\
9. Petra, Jordan
\annotation{Sung}\\
Carved in the rock where the desert winds blow,\\
An ancient kingdom from a long time ago.\\
Gotta see it now while the borders stay clear,\\
A world heritage that we hold so dear.

\annotation{Spoken}\\
8. Outback, Australia
\annotation{Sung}\\
Arid interior, wild and deep,\\
Where the survival documentaries sleep.\\
North to South or East to West,\\
Uluru’s Red Rock puts the soul to the test.

\annotation{Spoken}\\
7. Yellowstone National Park, USA
\annotation{Sung}\\
Oldest in the world, where the wild things grow,\\
Geological wonders and a winter snow.\\
Childhood dreams from a library book,\\
Grizzly bears waiting if you take a look.

\annotation{Spoken}\\
6. Patagonia, Argentina \& Chile
\annotation{Sung}\\
Southern Andes where the flat steppes meet,\\
Borders of the ocean where the currents compete.\\
By land through the woods or a cruise through the cold,\\
Glacier landscapes waiting to unfold.

\annotation{Spoken}\\
5. Taj Mahal, India
\annotation{Sung}\\
White marble tribute to a love so deep,\\
In the Golden Triangle where the ancient kings sleep.\\
Agra, Jaipur, Delhi in line,\\
An architectural wonder frozen in time.

\annotation{Spoken}\\
4. Angkor Wat, Cambodia
\annotation{Sung}\\
The biggest temple that the world has seen,\\
Reclaimed by the roots and the jungle green.\\
Stone-carved faces watching through the trees,\\
Moving like a ghost on a tropical breeze.

\annotation{Spoken}\\
3. Trans-Siberian Railway, Russia
\annotation{Sung}\\
Nine-thousand kilometers, Europe to the East,\\
A seven-day journey on a steel-wheeled beast.\\
From Moscow’s lights to the Japan Sea coast,\\
The longest active rhythm that the world can boast.

\annotation{Spoken}\\
2. Victoria Falls, Zambia \& Zimbabwe
\annotation{Sung}\\
The smoke that thunders, the water that falls,\\
Africa’s heartbeat echoing through the walls.\\
Zambezi river rushing down to the sea,\\
A misty white column of pure energy.

\annotation{Build-up: Snare roll intensifies}
\annotation{Spoken}\\
1. Caño Cristales, Colombia
\annotation{Sung}\\
Liquid rainbow in the ancient stone,\\
A crystal-clear river with a style of its own.\\
Red aquatic plants blooming under the sun,\\
Take a boat, take a flight, now the journey is done.
\annotation{The Drop}

\annotation{The Break: Honorable Mentions}
\annotation{The beat strips back to a minimal, pulsing synth loop as these names are quickly chanted in rhythm}

Cliffs of Dover (UK)\\
Bora Bora (Polynesia)\\
Galápagos (Ecuador)\\
Hawái (USA)\\
Lake Baikal (Russia)\\
Plitvice Lakes (Croatia)\\
Legoland (Denmark)\\
Machu Picchu (Peru)\\
Pyramids of Giza (Egypt)\\
Serengeti (Tanzania \& Kenya)...

\songpart{Outro: Fade out with the kick drum}\end{lyrics}

\section{Diamond Depth August}
\tag{original-work}\tag{salsa}

\subsection{Notes}
Original work.
Salsa composition with historical narrative framing.

\subsection{Style}
salsa dura, 104 BPM, percussion-led groove, timbales cascara, conga tumbao, bongos martillo, upright bass montuno, piano guajeos, brass stab chorus, call and response, closed-mic male vocals, live room bleed, tape saturation, spring reverb, syncopated clave, streetwise swagger

\begin{lyrics}
\songpart{Intro}
\annotation{instrumental brass attack}

\annotation{Spoken}\\
Three warriors opened up your eyes\\
to sword, to cross, to banner’s rise—\\
since then no fear your borders haunt,\\
nor greed your great heart’s noble want.

\songpart{Verse 1}
A diamond’s depth struck, one August,\\
the strings of a new lute’s trust,\\
and now we hear the melodious flow\\
in hymns where youth’s fresh voices go.

Fertile mother of proud offspring\\
who smile at vain and gilded thing—\\
still watchful for the morning light,\\
and to the past, forever bright.

\songpart{Chorus}
Let us sing a hymn to your sky,\\
to your land and your pure living—\\
white star that in the Andes glows!\\
broad path that to the future goes!

\songpart{Verse 2}
The savanna is a fallen sky,\\
a carpet spread where your feet lie;\\
and of the varied world you cheer,\\
you are both arm and brain, my dear.

Surviving from a golden reign,\\
a sunless empire’s ageless plane—\\
in you: a temple, shield, a grille,\\
a reredos, font, a lantern still.

\songpart{Chorus}
Let us sing a hymn to your sky,\\
to your land and your pure living—\\
white star that in the Andes glows!\\
broad path that to the future goes!

\songpart{Verse 3}
To Caldas, scanning stars on high,\\
to Bolívar, reborn nearby,\\
to Nariño, the press’s spark—\\
as in dreams, you see them in the dark.

Caros, Cuervos, Pombos, Silvas—\\
so many names of deathless fame;\\
on history’s endless, thread-like line,\\
your mother’s love gave life divine.

\songpart{Chorus}
Let us sing a hymn to your sky,\\
to your land and your pure living—\\
white star that in the Andes glows!\\
broad path that to the future goes!

\annotation{Montuno}
\annotation{Choir} Where is the melodious flow heard?
\annotation{Solo} In the hymns of youth, each word.
\annotation{Choir} How is the melodious flow heard?
\annotation{Solo} With the strings of a new lute stirred.

\annotation{Choir} What names are of immortal fame?
\annotation{Solo} Caros, Cuervos, Pombos—same.
\annotation{Choir} Is the savanna a fallen sky?
\annotation{Solo} A carpet spread for your feet, yes—why?

\songpart{Bridge}
Oriflamme of Gran Colombia—\\
in Caracas and Quito, yeah—\\
forever in your glory’s light,\\
with freedom’s reveilles burning bright.

\songpart{Verse 4}
Noble, loyal in peace and war,\\
at the foot of your strong hills’ lore,\\
and in the half-moon’s arched decree,\\
you rise as Cid, Santa Fe—free.

Flower of races, crown and sum—\\
in the homeland none else will come;\\
our voice the centuries repeat:\\
¡Bogotá! ¡Bogotá! ¡Bogotá!

\songpart{Chorus}
Let us sing a hymn to your sky,\\
to your land and your pure living—\\
white star that in the Andes glows!\\
broad path that to the future goes!

\songpart{Outro}
Flower of races, crown and sum—\\
in the homeland none else will come;\\
our voice the centuries repeat:\\
Bogotá! Bogotá! Bogotá!
\end{lyrics}

\section{Miladi in the Dark}
\tag{original-work}\tag{medieval-folk}

\subsection{Notes}
Original work.
Medieval folk ballad exploring grief, ritual, and dark-age emotional landscape.

\subsection{Style}
Medieval folk ballad with 6/8 sway, frame-story verses on lute, bowed fiddle, hand drum, and drone; verse 1 opens sparse and grave, pre-chorus adds psaltery and low harmonies, chorus lifts with a solemn canticle refrain and open vowel chant. Lead vocal is intimate and cracked, with doubled endings and soft group responses. Harp glints, wind chimes, and drum-roll transitions; warm, torchlit, earthy mix., medieval, emotional, ballad, sleep

\begin{lyrics}
\songpart{Verse 1}
In the dark she wept alone\\
Over him she had struck down\\
She cast her clothes aside\\
And clothed her skin in red\\
She lay atop his body\\
Kissing, shaking, praying\\
Her tears fell cold as ash\\
In the black, in the black

\songpart{Pre-Chorus}
Two years before that night\\
I rode to summer stone\\
To seek an old war friend\\
And plan our hunt at dawn\\
But there at his high hall\\
A girl stood in my sight\\
Seventeen springs of breath\\
And my heart lost its way

\songpart{Chorus}
Miladi, Miladi\\
White horse in the dawn\\
Miladi, Miladi\\
You looked me down and drawn\\
Miladi, Miladi\\
Your eyes were sharp as pine\\
Miladi, Miladi\\
And you were mine

\songpart{Verse 2}
I asked him of the girl\\
He frowned, then spoke her name\\
His daughter, born in war\\
Of Kamin far in Ismonia\\
The next day she rode out\\
With a crossbow at her side\\
White horse, white hands, no words\\
Just glances, just glances

\songpart{Pre-Chorus}
Twelve of us heard the horn\\
Then all the hounds broke loose\\
She spurred ahead of us\\
And I went hard behind\\
The others fell away\\
Till only woods remained\\
Their voices thin and far\\
Like smoke above the field

\songpart{Chorus}
Miladi, Miladi\\
White horse in the dawn\\
Miladi, Miladi\\
You looked me down and drawn\\
Miladi, Miladi\\
Your eyes were sharp as pine\\
Miladi, Miladi\\
And you were mine

\songpart{Bridge}
She stopped beside a log\\
And loosed her bolt to wood\\
Then loaded once again\\
And leveled steel at me
"Come closer," said her mouth
"Come closer and obey\\
Come kiss me, Grispo de Julminadi"\\
So I bent to her will

\songpart{Verse 3}
We kissed among the leaves\\
Till footsteps came near by\\
Then split apart at once\\
And hid our burning mouths\\
That evening after meat\\
I climbed alone to bed\\
Yet after hours had fled\\
She came and asked my company

\songpart{Final Chorus}
Miladi, Miladi\\
The dark knew both our names\\
Miladi, Miladi\\
We fed on secret flame\\
Miladi, Miladi\\
Then came the blade and night\\
Miladi, Miladi\\
And love lay down in blood
\end{lyrics}

\section{Destinies Overlap}
\tag{original-work}\tag{alternative-rock}

\subsection{Notes}
Original work.
Alternative rock narrative of parallel lives and unforeseen convergence.

\subsection{Style}
Alternative rock with a mid-tempo groove. Acoustic guitar strumming open chords in eighth notes. Electric bass guitar playing melodic lines with a warm, rounded tone. Drum kit featuring a dry snare and prominent hi-hat eighth notes. Male baritone vocals with a clean, dry production style. Electric guitar with light overdrive playing sustained lead lines. Key of G major. Tempo 134 BPM. 4/4 time signature. The arrangement builds from acoustic guitar and vocals to a full band texture.

\begin{lyrics}
\songpart{Chapter 1: The Hangover}
A drunken man awakes in grief and pain,\\
While elsewhere poise and purpose build a wall.\\
Two lives begin to lean toward their fall,\\
Divided by the miles and pouring rain.
\songpart{Chapter 2: Routine Breaking}
A cycle breaks as buried truths arise,\\
The past demands to be acknowledged now.\\
A solemn, silent, and unspoken vow,\\
Brings secrets out beneath the open skies.
\songpart{Chapter 3: Take-offs}
The journey starts with focus and with drive,\\
A mission takes its flight to distant lands.\\
The web is spun by heavy, hidden hands,\\
To keep the fragile hope of truth alive.
\songpart{Chapter 4: All Roads}
The paths converge as strangers draw so near,\\
Though unaware of who is watching them.\\
The tide begins to rise toward the hem,\\
Of all the hidden things they hold so dear.
\songpart{Chapter 5: Brunch}
The tensions hide behind a social mask,\\
A family gathers for a hollow meal.\\
With wounds that time has failed to truly heal,\\
They perform their part and finish every task.
\songpart{Chapter 6: Flashbacks}
The center point is found in blood and glass,\\
The moment that the narrative defines.\\
The truth is etched in sharp and bitter lines,\\
While shadows of the haunting memories pass.
\songpart{Chapter 7: A Pause}
A moment taken just to catch a breath,\\
The digital surveillance finds its pace.\\
A lonely study of a distant face,\\
That waits amidst the silence and the death.
\songpart{Chapter 8: High Stakes}
The gamble grows as dangers start to loom,\\
With secrets lost and sobriety gone.\\
They move across the dark and fading dawn,\\
To face the inevitable sense of doom.
\songpart{Chapter 9: Intermezzo}
The pace slows down within the quiet space,\\
Between the acts of what is soon to come.\\
The beating heart is steady as a drum,\\
While waiting for the end of this long race.
\songpart{Chapter 10: The Heat}
The lowland swelter hides a dangerous play,\\
Where fragments of the truth are caught by ear.\\
The moral weight is heavy and austere,\\
While logic slowly crumbles in the day.
\songpart{Chapter 11: Bon Soir}
A touch of grace within the city light,\\
The distances begin to softly blur.\\
The quiet intimacy starts to stir,\\
As shadows deepen in the coming night.
\songpart{Chapter 12: Vertrek}
A bittersweet departure marks the end,\\
Of connection forged in transient time.\\
A final note within a hollow chime,\\
As paths once joined now start to slowly bend.
\songpart{Chapter 13: The Glades}
The predators are circling in the deep,\\
With schemes that span across the ocean wide.\\
There is nowhere for guilt to run or hide,\\
When promises are ones they mean to keep.
 \songpart{Chapter 14: Agent Perez}
The institution stares with cold intent,\\
To measure out the cost of every crime.\\
A reckoning that waits for proper time,\\
With power that the hidden hands have lent.
\songpart{Chapter 15: Game On}
The trap is set as three sides pull the line,\\
The bait is taken with a single word.\\
The hidden calls and movements now are heard,\\
To seal the outcome of their dark design.
\songpart{Chapter 16: Morocho}
The puppet master waits within the pool,\\
The surface hides the depths of what he planned.\\
The power shifts within his steady hand,\\
As others play the part of every fool.
\songpart{Chapter 17: Los Diablos}
The tragedy is raw and stripped of pride,\\
The roles are played as darkness takes its toll.\\
The fractured pieces of a weary soul,\\
Have nowhere left that they can truly hide.
\songpart{Chapter 18: Roots}
The ancestors are calling from the past,\\
To ground the wanderer in soil and stone.\\
The seeds of all the lives they have now grown,\\
Are linked to where their roots were planted fast.
\songpart{Chapter 19: Movements}
The frantic pace begins to reach its peak,\\
As simultaneous ends start to arrive.\\
It is a desperate way to stay alive,\\
When futures are the only things they seek.
\songpart{Chapter 20: Operations}
The plans are turned to action and to force,\\
The blueprint moves to strike with sudden aim.\\
They play the final round within the game,\\
To steer the vessel on its deadly course.
\songpart{Chapter 21: Closings}
The doors are shut upon the lives they led,\\
The quiet end of all the work they did.\\
With all the secrets safely turned and hid,\\
They leave the ghosts of all the living dead.
\songpart{Chapter 22: The Stand}
The rope remains unpulled between the two,\\
They stand in silence waiting for the light.\\
The end is neither wrong and neither right,\\
As all the old is washed away by new.
\end{lyrics}

\section{Best Available Lens}
\tag{original-work}\tag{dubstep}

\subsection{Notes}
Original work.
Dubstep hymn on epistemology and practical philosophy.

\subsection{Style}
Dubstep hymn with half-time stomp, snarling wobble bass, and handclap-driven verses; verse rides sparse organ stabs and clipped spoken-sung lines, pre-chorus strips to sub pulses and rising choir pads, chorus hits with a huge detuned drop and stacked congregational vocals. Filtered risers, reverse cymbal swells, and glitchy vocal chops bridge the sections. Bright, विशाल, and sermon-like in the mix, with close-mic lead and wide, thunderous lows., dubstep, world, contemporary, theme

\begin{lyrics}
\songpart{Verse 1}
You keep calling it a creed\\
But it’s a working frame\\
A way to test the world\\
Not a holy name\\
It has rules for doing\\
What it sets out to do\\
Build a better map\\
Of rock and rain and you

\songpart{Pre-Chorus}
Philosophers ask why\\
And philosophers ask whether\\
That’s their long debate\\
They hold it all together\\
Scientists just work\\
By a learned, living line\\
Passed from hand to hand\\
Across the edge of time

\songpart{Chorus}
Best available lens\\
Best available lens\\
Not the final word\\
Still the way we learn\\
Best available lens\\
Best available lens\\
Useful, tested, made\\
For the world as it turns

\songpart{Verse 2}
Don’t wrap it in a banner\\
Don’t make it some new church\\
A methodology\\
Is not a sacred perch\\
Neopositivism\\
Pragmatism too\\
History left its fingerprints\\
In all we say and do

\songpart{Pre-Chorus}
And that shiny little brand\\
That called itself a thunder\\
Wasn’t a doctrine\\
Just a label dressed as wonder\\
New atheism fades\\
When the slogans lose their spark\\
A package, not a philosophy\\
No torch, no ark

\songpart{Chorus}
Best available lens\\
Best available lens\\
Not the final word\\
Still the way we learn\\
Best available lens\\
Best available lens\\
Useful, tested, made\\
For the world as it turns

\songpart{Bridge}
There are saints and skeptics\\
In the same white coat\\
Christian, Jew, and Buddhist\\
All of them afloat\\
Some chase rigor hard\\
Some chase the subject first\\
Some bring sloppy hands\\
And some bring a sharper thirst\\
Laypeople can twist it\\
Till it fits their tune\\
Truth when it flatters them\\
Then doubt beneath the moon\\
But a work in progress\\
Never asks to be divine\\
It asks to be improved\\
By the next in line

\songpart{Breakdown}
Not a religion\\
Not an ideology\\
An enterprise with a purpose\\
A method with a key\\
Not ultimate truth\\
But the best we’ve got\\
To learn what’s there\\
And what is not

\songpart{Final Chorus}
Best available lens\\
Best available lens\\
Not the final word\\
Still the way we learn\\
Best available lens\\
Best available lens\\
Useful, tested, made\\
For the world as it turns\\
Best available lens\\
Best available lens\\
Take the better map\\
Over every false concern
\end{lyrics}

\section{Not In The Creed}
\tag{original-work}\tag{gospel}\tag{soul}

\subsection{Notes}
Original work.
Gospel-soul reflection on religious childhood and eventual departure.

\subsection{Style}
soul-folk, gospel soul, 118 BPM, velvet bass baritone, gospel choir refrains, fingerpicked electric guitar, bowed fiddle lines, conga pulse, upright bass, brushed snare, organ swells, plate reverb, tape saturation, sidechain pumping, call and response, spoken intro, half-time bridge, bittersweet confession, devotional ache

\begin{lyrics}
\songpart{Verse 1}
\annotation{male vocals, velvet bass registry, soul delivery, smooth phrases}\\
Born where the bells rang hard\\
in a Catholic crowd, half-heart, half-card.\\
Christian Brothers taught my name—\\
first Communion, same old frame.\\
By fourteen, I’d made my mind\\
before Confirmation, I’d crossed that line.

\songpart{Pre-Chorus}
\annotation{male vocals, bass registry, soul delivery}\\
Still I stood there straight\\
Still I tried to wait\\
For some fire to land\\
In my open hand

\songpart{Chorus}
\annotation{gospel choir}\\
I was never all in\\
I was never all in\\
I wore the coat\\
But I missed the skin\\
I was never all in\\
I was never all in
(never all in)\\
Just telling the truth\\
About where I’ve been

\songpart{Verse 2}
\annotation{male vocals, velvet bass registry, soul delivery, smooth phrases}\\
Then Sweden took me in—\\
cold air, clean streets, less talk of sin.\\
Catholic turned to a sign—\\
more a name than a line.\\
Liberation books in my hands,\\
Cardijn’s kids, and their marching plans.\\
Back in Colombia, Jesuit halls—\\
still I heard a bigger call.

\songpart{Pre-Chorus}
\annotation{male vocals, bass registry, soul delivery}\\
Not the creed, not the show,\\
but the way we go.\\
Toward the hungry face.\\
Toward a harder grace.

\songpart{Chorus}
\annotation{golspel choir}\\
I was never all in\\
I was never all in\\
I wore the coat\\
But I missed the skin\\
I was never all in\\
I was never all in
(never all in)\\
Just telling the truth\\
About where I’ve been

\songpart{Bridge}
\annotation{rap, male baritone vocals}\\
I watched the online saints\\
The hard old rules, the holyaints\\
Faith in a fist\\
And a book like a list\\
They spoke of Jesus\\
Then played the Pharisee game\\
Fearing science\\
While blessing the shame\\
And I learned my fight\\
Wasn’t left or right\\
It was human first\\
It was daylight

\songpart{Verse 3}
\annotation{male vocals, velvet bass registry, soul delivery, smooth phrases}\\
If my branch was close to Him\\
Closer than some who loved the hymns\\
Still I couldn’t kneel\\
To what I couldn’t feel\\
No ghost in the room\\
No heaven in sight\\
Just conscience and work\\
And a harder light

\songpart{Pre-Chorus}
\annotation{male vocals, bass registry, soul delivery}\\
So I laid down the mask\\
Stopped the old half-ask\\
Called my own soul read\\
And said what I meant

\songpart{Chorus}
\annotation{Gosple choir}\\
I was never all in\\
I was never all in\\
I wore the coat\\
But I missed the skin\\
I was never all in\\
I was never all in
(never all in)\\
Just telling the truth\\
About where I’ve been

\songpart{Outro}
\annotation{male vocals, velvet bass registry}\\
Not in anger, not in spite\\
Just a cleaner name for night\\
I don’t call myself a Christian now\\
I call it honest, and I bow to how\\
I was never all in\\
I was never all in
(never all in)\\
Just telling the truth\\
About where I’ve been
\end{lyrics}

\section{Which God Stands Still}
\tag{original-work}\tag{gospel}

\subsection{Notes}
Original work.
Gospel composition on faith and the search for immovable truth.

\subsection{Style}
Gospel with bass baritone lead and responsive choir; slow-march 6/8 swing, handclaps on the backbeat, organ swells under sparse verse, pre-chorus opens with clapping and rising choir pads, chorus blooms with full SATB harmony and shouted responses, bridge drops to bass voice and room tone before a final lift. Close-mic lead in verses, gritty doubles on key lines, short delay throws on the hook, tambourine shakes, ascending organ runs, wide warm mix with a church-room bloom., worship, theme, strong, gospel

\begin{lyrics}
\songpart{Verse 1}
Which god wants fear?\\
Which god wants praise?\\
I read the stone\\
I read the page

No graven face\\
No golden frame\\
Just the hard old law\\
Calling my name

\songpart{Pre-Chorus}
You said, "Do not kneel"\\
To empty wood\\
You said, "Do not trade\\
The whole for the false"

If love is truth\\
Then love should lead\\
Not clutch my soul\\
And tighten its grasp

\songpart{Chorus}
Which God stands still?\\
Which God lets me grow?\\
Which God loves me\\
And lets me go? (lets me go)

Not a fist, not a throne\\
Not a child on a throne\\
Which God stands still?\\
Which God lets me grow? (lets me grow)

\songpart{Verse 2}
My catechism days\\
Had a cold bright seat\\
The rules were clear\\
The room was deep

The first commandment\\
Stood like a gate\\
No other gods\\
No idol bait

So I learned the fear\\
Then I learned the test\\
A love that demands\\
Can still be less

\songpart{Pre-Chorus}
If I lose my way\\
I can lose my name\\
But if I lose my mind\\
That is not grace

A loving parent\\
Would want my height\\
Would want my hands\\
To build in light

\songpart{Chorus}
Which God stands still?\\
Which God lets me grow?\\
Which God loves me\\
And lets me go? (lets me go)

Not a fist, not a throne\\
Not a child on a throne\\
Which God stands still?\\
Which God lets me grow? (lets me grow)

\songpart{Bridge}
\annotation{Drop to bass voice}\\
If there is One\\
Who made this breath\\
Then let it mean\\
More than dread

Let me reach\\
My full created flame\\
Let me love\\
Without a chain

\annotation{Choir swell}\\
Not every voice\\
That calls itself holy\\
Knows the shape\\
Of mercy

If You are Love\\
Then prove it slow\\
Make me whole\\
And let me grow

\songpart{Final Chorus}
Which God stands still?\\
Which God lets me grow?\\
Which God loves me\\
And lets me go? (lets me go)

Not a fist, not a throne\\
Not a child on a throne\\
Which God stands still?\\
Which God lets me grow? (lets me grow)

Which God stands still?\\
Which God lets me grow?\\
Amen, I’m asking\\
That I may know (that I may know)
\end{lyrics}

\section{Ashes in the Pews}
\tag{original-work}\tag{dubstep}\tag{gospel}

\subsection{Notes}
Original work.
Dubstep-gospel exploration of ritual transmission and inherited belief.

\subsection{Style}
Dubstep hymn at 140 BPM with half-time drops and swung handclap swing, verse rides sparse organ stabs and sub pulses, pre-chorus strips to soprano lead and low choir hum, chorus hits with stacked gospel choir, wobble bass, and clap-accented lifts. Soprano lead gets wide doubles, choir shouts on hook words, delay throws on the last line. Ear candy: reversed choir swell into each drop, chime flickers between phrases, seismic riser before final chorus. Bright, विशाल, and cathedral-wide, theme, gospel, world, worship, magical, dubstep

\begin{lyrics}
\songpart{Verse 1}
What makes people kneel?\\
Not the sky alone\\
Hands get taught to fold\\
By stories passed like bone

Not born with a temple\\
Not born with a throne\\
The first seed is planted\\
Then it starts to grow

\songpart{Pre-Chorus}
Show me the crack\\
Show me the seam\\
A world full of facts\\
Can wake a dream

\songpart{Chorus}
Who taught the gods?\\
Who taught the fear?
(Who taught the gods?)\\
Who put it here?

Scratch the paint off\\
See the room
(See the room)\\
Truth can bloom

\songpart{Verse 2}
Some break from the script\\
When the edges don't fit\\
When the old words buckle\\
And the page won't sit

When the promised fire\\
Looks like smoke and dust\\
When the holy map\\
Leads nowhere for us

\songpart{Pre-Chorus}
Show me the crack\\
Show me the seam\\
A world full of facts\\
Can wake a dream

\songpart{Chorus}
Who taught the gods?\\
Who taught the fear?
(Who taught the gods?)\\
Who put it here?

Scratch the paint off\\
See the room
(See the room)\\
Truth can bloom

\songpart{Bridge}
No, it's not a curse\\
Not a broken mind\\
Not a wound you can name\\
Just a turn of time

Some fear the void\\
Some fear the flood\\
Some worship the ending\\
Instead of the mud

\songpart{Final Chorus}
Who taught the gods?\\
Who taught the fear?
(Who taught the gods?)\\
Who put it here?

Scratch the paint off\\
See the room
(See the room)\\
Truth can bloom

Who taught the gods?\\
Who taught the fear?
(Who taught the gods?)\\
Why cure the clear?
\end{lyrics}

\section{Neon Shutter // Only Herself}
\tag*{2026}\tag{ai-portrait}
\subsection{Notes}
\subsection{Style}
Dark cinematic trip-hop ballad with a slow half-time pulse, clipped kick and brushed snare, sub-bass murmurs, and a glassy, looping synth motif that glows like shifting neon; verses stay intimate and close-mic’d with hushed, confessional female vocals, while choruses bloom into wider, reverb-laden harmonies. Filtered pads and distant vocal chops swell between sections, threaded with metallic ticks, rain-slick ambience, and low guitar swells for added emotional weight. The production moves from stark, shadowy minimalism into a cathartic, luminous final pass that feels like walking out of a dark alley into rain-washed city light — cool, nocturnal, and emotionally redemptive.

\begin{lyrics}
\songpart{Verse 1}
Neon on the metal, dripping down in pink and blue\\
Wet street in the silence, every echo feels like truth\\
Leather at my shoulders, collarbone a secret word\\
City staring through me, every siren going unheard

\songpart{Pre-Chorus}
I used to bend my edges just to fit inside their frame\\
Let them paint their colors over every bit of shame\\
Now the shutter’s rusted, but it’s shining just the same\\
I am not a canvas for a stranger’s little claims

\songpart{Chorus}
So I stand here in the half-light, split in blue and rose\\
Every broken version of me in these city glows\\
You don’t get to name me, I am not your hurt to own\\
I have walked through ashes just to call this steel my home\\
Look into my dark eyes, see the storm I made my friend\\
This is not the breakdown, this is where I start to mend

\songpart{Verse 2}
Kohl like tired comets, tracing nights I can’t forget\\
Hands that once were trembling, learning how to hold regret\\
Choker at my breathing, tiny rings that catch the rain\\
Every mark of metal is a map out of the pain

\songpart{Pre-Chorus}
I carved out a language where they only left me scars\\
Tiny burning letters hiding underneath the stars\\
All the times they told me I was just a fallen part\\
I kept one small ember pressed against my beating heart

\songpart{Chorus}
So I stand here in the half-light, split in blue and rose\\
Every broken version of me in these city glows\\
You don’t get to name me, I am not your hurt to own\\
I have walked through ashes just to call this steel my home\\
Look into my dark eyes, see the storm I made my friend\\
This is not the breakdown, this is where I start to mend

\songpart{Bridge}
They wanted pretty, they wanted soft\\
They built their ladders on my loss\\
But I learned breathing in the smoke\\
That you can heal where you were broke

I am the shutter, I am the street\\
I am the rain beneath your feet\\
I am the script under my skin\\
I am the door I’m walking in

\songpart{Chorus}
So I stand here in the half-light, split in blue and rose\\
Every shattered moment in me making something whole\\
You don’t get to frame me, I am not your scene to own\\
I have danced with shadows just to call this night my home\\
Look into my dark eyes, they are brightening again\\
This is not the ending, this is where I learn to mend

\songpart{Outro}
Pink across my cheekbones, blue across my quiet rage\\
I’m the only author writing lines on this new page\\
City hums indifferent, but my heart is beating loud\\
I am not forgiving them, I’m just walking through the crowd\\
Every step a promise I don’t have to disappear\\
Standing in the neon, I am wholly, fiercely here
\end{lyrics}

\section{Copper In The Pines}
\tag*{2026}\tag{ai-portrait}
\subsection{Notes}
\subsection{Style}
Indie-folk ballad with 92 BPM fingerpicked acoustic guitar, brushed frame drum, upright bass swells, and glowing string pads; verse stays intimate and close-mic, pre-chorus opens with harmony lift and soft hand percussion, chorus blooms with stacked female vocals and a chantable hook, bridge drops to bare voice and guitar before a final lift. Ear candy: reversed breathy swells into the chorus, delicate bell tones between lines, rainstick shimmer in transitions. Mix is warm, wide, and cinematic yet earthy., soft, natural, romantic, simple, warm, tone, light

\begin{lyrics}
\songpart{Verse 1}
You stand in the hush of pines\\
Auburn fire in the dawn\\
Freckles lit by borrowed gold\\
Like the forest made you born

Linen moving in the wind\\
Cream and oat against the green\\
Your eyes drift just past my own\\
Like a thought I’ve never seen

\songpart{Pre-Chorus}
Mist is sliding through the trunks\\
Soft as prayers across the moss\\
And the morning finds your hands\\
Like a secret it can cross

\songpart{Chorus}
You belong to the trees\\
You belong to the light\\
Copper in the pines\\
Burning soft and bright\\
You belong to the trees\\
You belong to me\\
Copper in the pines\\
Where the wild things breathe

\songpart{Verse 2}
Loose gold ring against your skin\\
Catching bits of heaven through\\
Every strand that frames your face\\
Feels like sunlight learned you too

Ancient bark and fern and stone\\
Hold you like they know your name\\
You don’t stand inside this world\\
You and it are just the same

\songpart{Pre-Chorus}
Mist is sliding through the trunks\\
Soft as prayers across the moss\\
And the morning finds your hands\\
Like a secret it can cross

\songpart{Chorus}
You belong to the trees\\
You belong to the light\\
Copper in the pines\\
Burning soft and bright\\
You belong to the trees\\
You belong to me\\
Copper in the pines\\
Where the wild things breathe

\songpart{Bridge}
And if I looked too hard\\
I’d break the spell in two\\
So I let the branches keep\\
The quiet part of you\\
The ferns, the haze, the dawn\\
The amber in your hair\\
I swear the forest only shines\\
Because you’re already there

\songpart{Final Chorus}
You belong to the trees\\
You belong to the light\\
Copper in the pines\\
Burning soft and bright\\
You belong to the trees\\
You belong to me\\
Copper in the pines\\
Where the wild things breathe
\end{lyrics}

\section{Peeling Paint Girl}
\tag*{2026}\tag{ai-portrait}
\subsection{Notes}
\subsection{Style}
Alt-rock with a slow swagger, clipped drum groove, rubbery bass, and jagged muted guitars; verse stays close-mic and tense, pre-chorus opens with rising toms and breathy doubles, chorus hits with gang vocals and a rough octave harmony on the hook. Use scraped guitar harmonics, a reverse cymbal into each chorus, and a dusty room slap on the lead. Mix is gritty, wide, and filmic with sharp vocal presence., distinct, soft, raw, moody, low, natural

\begin{lyrics}
\songpart{Verse 1}
Pale skin, hard stare\\
Green-gray eyes like knives\\
Grease in the black hair\\
Damp in the evening light

Silver hoops at her ears\\
Distressed leather, torn at the seam\\
She stands by the old wall\\
Like she knows what it means

\songpart{Pre-Chorus}
Weathered paint on concrete\\
Peeling down in strips\\
Long shadows on her collar\\
Tension in her lips

She don’t flinch\\
She don’t fold\\
Something in her gaze\\
Stays cold

\songpart{Chorus}
Hold that look, hold that look\\
Don’t let it slip away\\
Hold that look, hold that look\\
Rough and bright and gray\\
She wears the scarred-up jacket\\
Like a second skin\\
Hold that look, hold that look\\
Let the tension in

\songpart{Verse 2}
Beige tank, raw-edge cut\\
Nothing dressed to please\\
Slim frame, steady weight\\
Standing in the breeze

Every scratch on the wall\\
Feels like it has a name\\
She’s got that worn-down fire\\
Burning just the same

\songpart{Pre-Chorus}
Weathered paint on concrete\\
Peeling down in strips\\
Long shadows on her collar\\
Tension in her lips

She don’t flinch\\
She don’t fold\\
Something in her gaze\\
Stays cold

\songpart{Chorus}
Hold that look, hold that look\\
Don’t let it slip away\\
Hold that look, hold that look\\
Rough and bright and gray\\
She wears the scarred-up jacket\\
Like a second skin\\
Hold that look, hold that look\\
Let the tension in

\songpart{Bridge}
Golden light on busted walls\\
Soft on all that grit\\
Beauty with the edges on\\
That’s the part that hits

She’s all smoke and steel\\
All quiet thunder, too\\
And every line on her face\\
Feels brutally true

\songpart{Final Chorus}
Hold that look, hold that look\\
Don’t let it slip away\\
Hold that look, hold that look\\
Rough and bright and gray\\
She wears the scarred-up jacket\\
Like a second skin\\
Hold that look, hold that look\\
Let the tension in
\end{lyrics}

\section{Ulla In The Hall}
\tag{encounters}\tag{ulla}\tag*{2026}\tag{ai-portrait}
\subsection{Notes}
\subsection{Style}
Dubstep anthem with heavy half-time drops, wobbling sub-bass, clipped synth stabs, and velvet female vocals; verse rides sparse bass pulses and dry snaps, pre-chorus strips to filtered pads and breathy doubles, chorus hits with a giant wobble lead and gang-voiced hook repeats. Use whispered ad-libs, delay throws on the name, reverse swells into each drop, and a glossy close-mic mix that feels cinematic and sharp., beautiful, soft, playful, fresh, piercing, dubstep, deep, light
\begin{lyrics}
\songpart{Verse 1}
Right in the hall\\
Fresh from the shower\\
Soft glow on skin\\
Rosy and bright

Hair on your shoulders\\
Loose and still damp\\
That oversized hoodie\\
Barely behaves

Hazel eyes on me\\
Straight nose, that smirk\\
Little lip spark\\
Says you know

You stand like trouble\\
Wrapped in a dream\\
Still somehow cool\\
And totally free

\songpart{Pre-Chorus}
Hold that look\\
Don’t move\\
Let the room lean in\\
Let it all zoom

Slow that grin\\
One step\\
Feel the hush go still\\
Before you hit

\songpart{Chorus}
Ulla in the hall\\
Ulla in the hall\\
Fresh glow, play it cool\\
Ulla in the hall

Ulla in the hall\\
Ulla in the hall\\
Barely covered, still gold\\
Ulla in the hall

(Oh-oh, Ulla)
(Oh-oh, Ulla)

\songpart{Verse 2}
Fair skin shining\\
Like rain on glass\\
Slim and steady\\
Built to last

Close-up magic\\
Every detail cuts\\
Soft feminine lines\\
And fearless touch

You own that lobby\\
Like it knows your name\\
One tiny smile\\
And I’m caught

Playful and poised\\
With that quiet fire\\
You turn the whole frame\\
Into art

\songpart{Pre-Chorus}
Hold that look\\
Don’t blink\\
Let the bass come back\\
Let the floor sink

Slow that grin\\
One pulse\\
Feel the hush go still\\
Before it breaks

\songpart{Chorus}
Ulla in the hall\\
Ulla in the hall\\
Fresh glow, play it cool\\
Ulla in the hall

Ulla in the hall\\
Ulla in the hall\\
Barely covered, still gold\\
Ulla in the hall

(Oh-oh, Ulla)
(Oh-oh, Ulla)

\songpart{Bridge}
Close focus\\
No blur\\
Every line stays true\\
Every glance hits hard

Soft inside\\
But sharp\\
You don’t need to try\\
You already are

\songpart{Break}
(Ulla)
(Show me that smirk)
(Ulla)
(Shine in the hush)

\songpart{Final Chorus}
Ulla in the hall\\
Ulla in the hall\\
Fresh glow, play it cool\\
Ulla in the hall

Ulla in the hall\\
Ulla in the hall\\
Barely covered, still gold\\
Ulla in the hall

(Oh-oh, Ulla)
(Oh-oh, Ulla)
\end{lyrics}
